%%%%%%%%%%%%%%%%%%%%%%%%%%%%%%%%%%%%%%%%%%%%%%%%%%%%%%%%%%%%%%%%%%%%%%
% Lớp 11 | Các số đặc trưng đo xu thế trung tâm cho mẫu số liệu | Số trung bình và mốt của mẫu số liệu ghép nhóm
% Dạng 1D5?1-1: Câu hỏi lý thuyết
% Số câu: 87  (sinh tự động bởi tach_cau.py; từ nay thầy sửa trực tiếp file này)
%%%%%%%%%%%%%%%%%%%%%%%%%%%%%%%%%%%%%%%%%%%%%%%%%%%%%%%%%%%%%%%%%%%%%%

%%%%%%%%%%%%%%%% TRẮC NGHIỆM 4 PHƯƠNG ÁN %%%%%%%%%%%%%%%%

% nguồn: Ok_Toan11_TracNghiem_Lan2
\begin{ex}%[1D5N1-1]%[ma:52da42ab]
	Một nhà thực vật học đo chiều dài của $74$ chiếc lá và thu được kết quả như sau (đơn vị: mm) 
	\begin{center}
		\begin{tabular}{|l|c|c|c|c|c|c|c|}
			\hline Chiều dài (mm)  & {$[5{,}45;5{,}85)$} & {$ [5{,}85;6{,}25)$} & {$[6{,}25;6{,}65) $} & {$[6{,}65;7{,}05)$} & {$[7{,}05;7{,}45)$} & {$[7{,}45);7{,}85)$}\\
			\hline Tần số & $5$ & $9$ & $15$ &  $19$& $16$& $8$ \\
			\hline 
		\end{tabular}
	\end{center}
	Giá trị đại  diện và độ dài của nhóm $[6{,}25;6{,}65)$là
	 \choice
		{ $6{,}05$ và $0{,}2$} 
		 { $6{,}45$ và $0{,}2$}
		{ $6{,}05$ và $0{,}4$} 
		{  \True $6{,}45$ và $0{,}4$}
	 \loigiai{
		Giá trị đại diện của nhóm là  $\dfrac{6{,}25+6{,}65}{2}=6{,}45$.\\
	Độ dài  của nhóm là $6{,}65-6{,}25=0{,}4$.}
\end{ex}

% nguồn: Ok_Toan11_TracNghiem_Lan2
\begin{ex}%[1D5N1-1]%[ma:17b863db]
	Giả sử  nhóm $k$ là nhóm đầu tiên có tần số tích lũy lớn hơn hoặc bằng $\dfrac{n}{2}$, tức là $cf_{k-1} < \dfrac{n}{2}$ nhưng $cf_k \ge \dfrac{n}{2}$. Gọi $r$, $d$, $n_k$ lần lượt là đầu mút trái, độ dài, tần số của nhóm $k$; $cf_{k-1}$ là tần số tích lũy của nhóm $k-1$. Công thức trung vị nào sau đây là đúng?
	\choice 
		{\True $M_e=r+\left(\dfrac{\dfrac{n}{2}-cf_{k-1}}{n_k}  \right)\cdot d$}
		 {$M_e=r+\left(\dfrac{\dfrac{n}{2}+cf_{k-1}}{n_k} \right)\cdot d$}
		{$M_e=d+\left(\dfrac{\dfrac{n}{2}-cf_{k-1}}{n_k}  \right)\cdot r$}
		{$M_e=d+\left(\dfrac{\dfrac{n}{2}+cf_{k-1}}{n_k} \right)\cdot  r$}
	\loigiai{ 
	 }
\end{ex}

% nguồn: Ok_Toan11_TracNghiem_Lan2
\begin{ex}%[1D5N1-1]%[ma:1e153145]
	Quãng đường (km) từ nhà đến nơi làm việc của $40$ công nhân một nhà máy  được ghi lại như sau
	\begin{center}
		\begin{tabular}{|l|c|c|c|c|c|c|c|c|c|c|c|c|c|c|c|c|c|c|c|c|}
			 \hline $5$&$3$&$10$&$20$&$25$&$11$&$13$&$7$&$12$&$31$&$19$&$10$&$12$&$17$&$18$&$11$&$32$&$17$&$16$&$2$\\
			\hline $7$&$9$&$7$&$8$&$3$&$5$&$12$&$15$&$18$&$3$&$12$&$14$&$2$&$9$&$6$&$15$&$15$&$7$&$6$&$12$\\ 
			\hline 
		\end{tabular}
	\end{center}
	Ghép  nhóm số liệu trên thành các nhóm\\
	$[0;5),[5;10),[10;15),[15;20),[20;25),[25;30),[30;35)$. Tần số của  nhóm $[15;20)$ là
	 \choice
		 { $8$}
		 { $6$}
		 { $10$}
		 { \True $9$}
	\loigiai{ 
	Tần số  của nhóm $[15;20)$ là $n=9$.}
\end{ex}

% nguồn: Ok_Toan11_TracNghiem_Lan2
\begin{ex}%[1D5N1-1]%[ma:f5b5d541]
	Kết quả kì thi trắc nghiệm môn Toán với thang điểm $100$  của $32$ học sinh được cho trong bảng sau
	\begin{center}
		\begin{tabular}{|l|c|c|c|c|c|c|}
			\hline Điểm & {$[40;50)$} & {$ [50;60)$}  & {$[60;70) $} & {$[70;80)$} & {$[80;90)$} & {$[90;100)$}\\
			\hline Tần số & $4$ & $6$ & $11$ &  $6$& $3$& $2$ \\
			 \hline
		\end{tabular}
	\end{center}
	Độ  dài của nhóm $[90;100)$ là
	\choice 
		 { $7$}
		{  $8$}
		{  $9$}
		{  \True $10$}
	\loigiai{ 
	Độ dài của nhóm $[90;100)$ là  $100-90=10$.}
\end{ex}

% nguồn: Ok_Toan11_TracNghiem_Lan2
\begin{ex}%[1D5N1-1]%[ma:19f0df06]
	Giả sử mẫu số liệu được cho dưới dạng bảng  tần số ghép nhóm
	\begin{center}
		\begin{tabular}{|c|c|c|c|c|}
			 \hline
			Nhóm & Nhóm 1 &  Nhóm 2 & $\ldots$ & Nhóm $k$ \\
			 \hline
			 Giá trị đại diện & $c_1$ & $c_2$ & $\ldots$ & $c_k$ \\
			\hline 
			Tần số & $n_1$ & $n_2$ & $\ldots$ & $n_k$  \\
			\hline 
		\end{tabular}
	\end{center}
	Đặt $n=n_1+n_2+\cdots+n_k$. 
	Số trung bình của mẫu số liệu ghép  nhóm, kí hiệu $\overline{x}$ được tính theo công thức nào
	\choice 
		{\True  $\overline{x}=\dfrac{n_1 c_1+n_2 c_2+\cdots+n_k c_k}{n}$}
		{$\overline{x}=\dfrac{n_1 c_1+n_2 c_2+\cdots+n_k  c_k}{2n}$}
		 {$\overline{x}=\dfrac{n_1^2c_1+n_2^2c_2+\cdots+n_k^2c_k}{n}$}
		{$\overline{x}=\dfrac{n_1 c_1+n_2 c_2+\cdots+n_k c_k}{\sqrt{n}}$} 
	\loigiai{ 
		Số trung bình của mẫu số liệu ghép nhóm, kí hiệu $\overline{x}$, được tính  như sau
		\[\overline{x}=\dfrac{n_1 c_1+n_2  c_2+\cdots+n_k c_k}{n}, \text{ trong đó } n=n_1+n_2+\cdots+n_k.\]
	} 
\end{ex}

% nguồn: Ok_Toan11_TracNghiem_Lan2
\begin{ex}%[1D5N1-1]%[ma:ab6f82ca]
	Tuổi thọ (năm) của $50$ bình ắc quy ô tô được cho như sau 
	\begin{center}
		\begin{tabular}{|l|c|c|c|c|c|c|c|}
			 \hline Tuổi thọ (năm) & {$[2;2{,}5)$} & {$ [2{,}5;3)$} & {$[3;3{,}5) $} & {$[3{,}5;4)$} & {$[4;4{,}5)$} & {$[4{,}5;5)$}\\
			\hline Tần số & $4$ &  $9$ & $14$ & $11$& $7$& $5$ \\
			 \hline
		\end{tabular}
	\end{center}
	Giá trị đại diện của nhóm  $[2,5;3)$ là
	\choice 
		 { $2{,}5$}
		{ \True  $2{,}75$}
		 { $9$}
		 { $3$}
	 \loigiai{
	 Giá trị đại diện của nhóm $[2{,}5;3)$ là $\dfrac{2{,}5+3}{2}=2{,}75$.}
\end{ex}

% nguồn: Ok_TracNghiemLop11_Lan1
\begin{ex}%[1D5N1-1]%[ma:b34d2fb9]
	Cho mẫu số  liệu ghép nhóm về số tiền mà sinh viên chi cho thanh toán cước điện thoại trong tháng
	\begin{center}
		\begin{tabular}{|c|c|c|c|c|c|}
			 \hline
			Số tiền (Nghìn đồng) & $[0{;}50$ & $[50{;}100)$ &  $[100{;}150)$ & $[150{;}200)$ & $[200{;}250)$ \\
			\hline 
			Số sinh viên & $5$ & $12$ & $23$ & $17$ & $31$ \\ 
			\hline 
		\end{tabular}
	\end{center}
	Có bao nhiêu sinh viên chi từ $100$ đến dưới  $150$ nghìn đồng cho việc thanh toán cước điện thoại trong tháng?
	\choice 
		{$5$} 
		{\True  $23$}
		{$12$} 
		{$17$} 
	\loigiai{ 
		Có $23$ sinh viên chi từ $100$ đến dưới $150$ nghìn đồng cho việc thanh toán  cước điện thoại trong tháng.
	} 
\end{ex}

% nguồn: Ok_TracNghiemLop11_Lan1
\begin{ex}%[1D5N1-1]%[ma:e98e8502]
	 Mỗi nhóm số liệu ghép nhóm là tập hợp gồm:
	\choice 
		{các giá trị  của số liệu được ghép nhóm theo nhiều tiêu chí xác định}
		{các giá trị của số liệu được  ghép nhóm theo hai tiêu chí xác định}
		{\True các giá trị của số liệu được ghép nhóm theo một  tiêu chí xác định}
		{các giá trị của số liệu được ghép  nhóm theo ba tiêu chí xác định}
	 \loigiai{
		Mỗi nhóm số liệu là tập hợp gồm các giá trị của số liệu được ghép nhóm theo một tiêu chí xác  định.
	} 
\end{ex}

% nguồn: Ok_TracNghiemLop11_Lan1
\begin{ex}%[1D5N1-1]%[ma:9a698735]
	Đo chiều cao (tính bằng cm)  của $500$ học sinh trong một trường THPT ta thu được kết quả như sau
	\begin{center}
		\begin{tabular}{|c|c|c|c|c|c|}
			\hline 
			Chiều cao &  $\left[150;154\right)$ & $\left[154;158\right)$ & $\left[158;162\right)$ & $\left[162;166\right)$ & $\left[166;170\right)$ \\
			 \hline
			Số học sinh & $25$ & $50$ & $200$ & $175$ & $50$ \\ 
			 \hline
		\end{tabular}
	\end{center}
	Mẫu số liệu  ghép nhóm đã cho có tất cả bao nhiêu nhóm?
	\choice 
		 {\True $5$}
		 {$6$}
		 {$7$}
		{$12$} 
	 \loigiai{
		Mẫu số liệu ghép nhóm đã cho  có tất cả $5$ nhóm.
	 }
\end{ex}

% nguồn: Ok_TracNghiemLop11_Lan1
\begin{ex}%[1D5N1-1]%[ma:908873e9]
	Điều tra về chiều cao của học sinh khối lớp 11, ta được mẫu số liệu sau: 
	\begin{center}
		\begin{tabular}{|c|c|}
			 \hline
			Chiều cao &  Số học sinh \\
			\hline 
			$[150{;}152$ & $5$  \\
			\hline 
			$[152{;}154$ & $18$  \\
			 \hline
			 $[154{;}156$ & $40$ \\
			 \hline
			$[156{;}158$ & $26$ \\ 
			 \hline
			$[158{;}160$ & $8$ \\ 
			 \hline
			 $[160{;}162$ & $3$ \\
			\hline 
		\end{tabular}
	\end{center}
	Mẫu  số liệu ghép nhóm đã cho có bao nhiêu nhóm?
	 \choice
		 {$5$}
		{\True  $6$}
		{$7$} 
		{$12$} 
	 \loigiai{
		Có $6$ nhóm số  liệu ở bảng trên là $[150{;}152$, $[152{;}154$, $[154{;}156$, $[156{;}158$, $[158{;}160$, $[160{;}162$.
	} 
\end{ex}

% nguồn: Ok_TracNghiemLop11_Lan1
\begin{ex}%[1D5N1-1]%[ma:36fad1e2]
	Giả sử mẫu số liệu được cho dưới dạng bảng  tần số ghép nhóm.
	\begin{center}
		 {\tabcolsep = 5mm
		\begin{tabular}{|c|c|c|c|c|}
			 \hline
			Nhóm & Nhóm $1$ & Nhóm $2$ & \ldots & Nhóm  $k$\\
			 \hline
			 Giá trị đại diện & $c_1$ & $c_2$ & \ldots & $c_k$\\
			 \hline
			Tần  số & $n_1$ & $n_2$ & \ldots & $n_k$\\
			\hline 
		\end{tabular}
		} 
	\end{center}
	Đặt $n=n_1+n_2+ \cdots  + n_k$.\\
	 Số trung bình của mẫu số liệu ghép nhóm, kí hiệu $\overline{x}$, được tính theo công thức nào?
	 \choice
		{\True $\overline{x}=\dfrac{n_1 c_1+n_2 c_2 +\cdots +  n_k c_k}{n}$}
		{$\overline{x}=\dfrac{n_1 c_1+n_2 c_2+\cdots + n_k  c_k}{2n}$}
		{$\overline{x}=\dfrac{n_1^2 c_1+n_2^2  c_2+\cdots + n_k^2 c_k}{n}$}
		{$\overline{x}=\dfrac{n_1 c_1+n_2  c_2 +\cdots + n_k c_k}{\sqrt n}$}
	 \loigiai{
		Giả sử mẫu số liệu được  cho dưới dạng bảng tần số ghép nhóm
		\begin{center}
			 {\tabcolsep = 5mm
			\begin{tabular}{|c|c|c|c|c|}
				 \hline
				Nhóm & Nhóm  $1$ & Nhóm $2$ & \ldots & Nhóm $k$\\
				\hline 
				Giá trị đại diện & $c_1$ & $c_2$ & \ldots  & $c_k$\\
				 \hline
				Tần số & $n_1$ & $n_2$ & \ldots & $n_k$\\ 
				\hline 
			\end{tabular}
			 }
		\end{center}
		Số trung bình của mẫu số liệu ghép nhóm,  kí hiệu $\overline{x}$, được tính như sau
		$$\overline{x}=\dfrac{n_1 c_1+n_2 c_2 +\cdots + n_k c_k}{n}$$ 
		trong đó  $n=n_1+n_2+ \cdots + n_k$.
	 }
\end{ex}

% nguồn: Ok_TracNghiemLop11_Lan1
\begin{ex}%[1D5N1-1]%[ma:9ec68524]
	Cho mẫu số liệu ghép nhóm về thời gian (phút) đi từ nhà đến nơi làm việc của các nhân viên một công  ty như sau:
	\begin{center}
		\begin{tabular}{|c|c|c|c|c|c|c|c|}
			 \hline
			Thời gian & $[15{;}20$ & $[20{;}25)$  & $[25{;}30)$ & $[30{;}35)$ & $[35{;}40)$ & $[40{;}45)$ & $[45{;}50)$ \\
			 \hline
			Số nhân viên & $6$ & $14$  & $25$ & $37$ & $21$ & $13$ & $9$ \\
			 \hline
		\end{tabular}
	\end{center}
	 Mẫu số liệu được chia thành bao nhiêu nhóm?
	\choice 
		{$6$  nhóm}
		{$5$  nhóm}
		 {\True $7$ nhóm}
		{$8$  nhóm}
	\loigiai{ 
		 Mẫu số liệu được chia thành $7$ nhóm.
	} 
\end{ex}

% nguồn: Ok_TracNghiemLop11_Lan1
\begin{ex}%[1D5N1-1]%[ma:2ce191ee]
	Dựa trên bảng số liệu về  chiều cao của $100$ học sinh một trường trung học phổ thông dưới đây.
	\begin{center}
		\begin{tabular}{|c|c|c|}
			\hline 
			Nhóm & Chiều cao  (cm) & số học sinh \\
			 \hline
			$1$ &  $[150;153)$ & $7$ \\
			 \hline
			$2$  & $[153;156)$ & $13$ \\
			 \hline
			 $3$ & $[156;159)$ & $40$ \\
			\hline 
			$4$ &  $[159;162)$ & $21$ \\
			 \hline
			 $5$ & $[162;165)$ & $13$ \\
			 \hline
			$6$ & $[165;168)$ & $6$  \\
			 \hline
		\end{tabular}
	\end{center}
	$160{,}5$ là giá trị đại diện cho  nhóm
	\choice 
		 {$2$}
		{$3$} 
		{\True  $4$}
		 {$5$}
	 \loigiai{
		$160{,}5$ là giá trị đại diện cho nhóm $4$ vì $160{,}5=\dfrac{159+162}{2}$. 
	} 
\end{ex}

% nguồn: Ok_TracNghiemLop11_Lan1
\begin{ex}%[1D5N1-1]%[ma:c22fc398]
	Mẫu số liệu (T) được mô tả dưới dạng bảng thống kê  sau:
	\begin{center}
		\begin{tabular}{|c|c|c|c|c|c|c|}
			\hline 
			Tổng  điểm & $<6$ & $[6{;}7)$ & $[7{;}8)$ & $\cdots$ & $[28{;}29)$ & $[29{;}30)$ \\
			 \hline
			Số học sinh & $23$ & $69$ &  $129$ & $\cdots$ $216$ & $12$ \\
			 \hline
		\end{tabular}
	\end{center}
	Tổng số giá trị của mẫu số liệu là: 
	\choice 
		 {$2589863$}
		{$14586632$} 
		 {$125863$}
		{\True $1379008$} 
	 \loigiai{
		Tổng số giá trị của mẫu số liệu là: $344752 \cdot 4  = 1379008$ giá trị.
	} 
\end{ex}

% nguồn: Ok_TracNghiemLop11_Lan1
\begin{ex}%[1D5N1-1]%[ma:5becb02e]
	Mẫu số liệu được mô tả dưới  dạng bảng thống kê sau:
	\begin{center}
		\begin{tabular}{|c|c|c|c|c|c|c|}
			\hline 
			Tổng  điểm & $<6$ & $[6{;}7)$ & $[7{;}8)$ & $\cdots$ & $[28{;}29)$ & $[29{;}30)$ \\
			 \hline
			Số học sinh & $23$ & $69$ & $129$ & $\cdots$ & $216$ & $12$  \\
			\hline 
		\end{tabular}
	\end{center}
	Số lượng thí sinh có ít nhất một  môn học có điểm từ $6$ đến dưới $7$ là:
	\choice 
		{$23$} 
		{$192$} 
		{$56$} 
		{\True  $69$}
	\loigiai{ 
		Số lượng thí sinh có  ít nhất một môn học có điểm từ $6$ đến dưới $7$ là $69$
	} 
\end{ex}

% nguồn: Ok_TracNghiemLop11_Lan1
\begin{ex}%[1D5N1-1]%[ma:aa18bb7e]
	Nhóm số liệu ghép nhóm thường được cho  dưới dạng
	\choice 
		 {$[a{;}b]$, trong đó $a$ là đầu mút trái, $b$ là đầu mút phải}
		{$(a{;}b)$, trong đó $a$ là đầu  mút trái, $b$ là đầu mút phải}
		{$(a{;}b]$, trong đó $a$ là đầu  mút trái, $b$ là đầu mút phải}
		{\True $[a{;}b)$, trong đó $a$ là  đầu mút trái, $b$ là đầu mút phải}
	\loigiai{ 
		Nhóm số liệu thường được cho dưới dạng $[a{;}b)$, trong đó $a$ là đầu mút trái, $b$ là đầu mút  phải.
	} 
\end{ex}

% nguồn: Ok_TracNghiemLop11_Lan1
\begin{ex}%[1D5N1-1]%[ma:cfadac67]
	Thời gian hoàn thành bài kiểm tra Toán $45$ phút của các bạn trong lớp  được cho như sau
	\begin{center}
		\begin{tabular}{|c|c|c|c|c|c|}
			 \hline
			Thời gian (phút) & $[25; 30)$ & {$[30; 35)$} & {$[35; 40)$} &  {$[40; 45]$} \\
			\hline 
			Số học sinh & $2$ & $7$ & $10$  & $25$ \\
			\hline 
		\end{tabular}
	\end{center}
	Có bao nhiêu học sinh hoàn thành bài kiểm tra trước  khi hết giờ trên 5 phút?
	 \choice
		 {\True $19$}
		 {$18$}
		 {$20$}
		{$17$} 
	\loigiai{ 
		\begin{enumerate}
			\item Các nhóm số liệu $[25;  30)$, $[30; 35)$, $[35; 40)$, $[40; 45]$ với tần số tương ứng là $2,7,10,25$.
			\item Số học sinh hoàn thành bài kiểm tra trước khi hết giờ ít nhất $5$ phút là  $2+7+10=19$.
		\end{enumerate}
	} 
\end{ex}

% nguồn: Ok_TracNghiemLop11_Lan1
\begin{ex}%[1D5N1-1]%[ma:a3960bac]
	Khảo sát thời gian tập thể dục của một số học sinh khối  $11$ thu được mẫu số liệu ghép nhóm sau:
	\begin{center}
		\begin{tabular}{|c|c|c|c|c|c|}
			\hline  Thời gian & {$[0; 20)$} & {$[20; 40)$} & {$[40; 60)$} & {$[60; 80)$} & {$[80;100)$} \\
			\hline Số học  sinh & $5$ & $9$ & $12$ & $10$ & $6$ \\
			\hline 
		\end{tabular}
	\end{center}
	Giá trị đại diện của nhóm  $\left[40\,;\,60\right)$ là
	\choice 
		{$0$} 
		 {$42$}
		{\True $50$} 
		 {$100$}
	 \loigiai{
	Giá trị đại diện của nhóm $\left[40\,;\,60\right)$  là $\dfrac{40+60}{2}=50$.}
\end{ex}

% nguồn: Ok_TracNghiemLop11_Lan1
\begin{ex}%[1D5N1-1]%[ma:a608fb42]
	Tìm hiểu thời gian xem tivi trong tuần trước (đơn  vị: giờ) của một số học sinh thu được kết quả sau:
	\begin{center}
		\begin{tabular}{|c|c|c|c|c|c|}
			\hline 
			Thời gian (giờ) & $[0{;}5$ & $[5{;}10)$ & $[10{;}15)$ &  $[15{;}20)$ & $[20{;}25)$ \\
			\hline 
			Số học sinh & $8$ &  $16$ & $4$ & $2$ & $2$ \\
			\hline 
		\end{tabular}
	\end{center}
	Có bao nhiêu học sinh có thời gian xem ti vi từ 20 giờ  đến dưới 25 giờ trong tuần trước?
	 \choice
		{$5$} 
		{$4$} 
		{$3$} 
		{\True $2$} 
	\loigiai{ 
		Từ bảng số liệu ta thấy có $2$ học sinh có thời gian xem ti vi từ $20$ giờ  đến dưới $25$ giờ trong tuần trước.
	 }
\end{ex}

% nguồn: Ok_TracNghiemLop11_Lan1
\begin{ex}%[1D5N1-1]%[ma:d4874262]
	Điền từ  thích hợp và chỗ trống: "Mẫu số liệu... là mẫu số liệu cho dưới dạng bảng tần số của các nhóm số liệu."
	\choice 
		{Không ghép  nhóm}
		{\True Ghép  nhóm}
		 {Ghép nhóm và không ghép nhóm}
		{Cả ba câu trên đều  sai}
	 \loigiai{
		 Mẫu số liệu ghép nhóm là mẫu số liệu cho dưới dạng bảng tẩn số của các nhóm số liệu.
	 }
\end{ex}

% nguồn: Ok_TracNghiemLop11_Lan1
\begin{ex}%[1D5N1-1]%[ma:7226d869]
	Một công ty may quần áo đồng phục học sinh cho biết  cỡ áo theo chiều cao của học sinh được tính như sau:
	\begin{center}
		\begin{tabular}{|c|c|c|c|c|c|}
			\hline 
			Chiều cao (cm) &  $[150{;}160$ & $[160{;}167)$ & $[167{;}170)$ & $[170{;}175)$ & $[175{;}180)$ \\
			\hline 
			Cỡ áo & $S$ & $M$ & $L$ & $XL$  & $XXL$ \\
			\hline 
		\end{tabular}
	\end{center}
	Công ty muốn ước lượng tỉ lệ các cỡ áo khi may cho học sinh lớp $11$ đã đo chiều cao của $36$ học  sinh khối $11$ của một trường và thu được mẫu số liệu sau (đơn vị là cm)
	\begin{center}
		\begin{tabular}{|c|c|c|c|c|c|c|c|c|}
			 \hline
			$160$ & $161$  & $161$ & $162$ & $162$ & $162$ & $163$ & $163$ & $163$ \\
			\hline 
			$164$ & $164$ & $164$ & $164$ & $165$ & $165$ & $165$ & $165$ & $165$ \\ 
			\hline 
			$166$ & $166$ & $166$ & $166$  & $167$ & $167$ & $168$ & $168$ & $168$ \\
			\hline 
			$168$ & $169$ & $169$ & $170$ & $171$  & $171$ & $172$ & $172$ & $174$ \\
			 \hline
		\end{tabular}
	\end{center}
	Số học sinh có  chiều cao vừa vặn với size L là bao nhiêu?
	\choice 
		{\True  $8$}
		{$7$} 
		 {$6$}
		 {$5$}
	\loigiai{ 
		Số học sinh có chiều cao vừa vặn với size $L$ là $8$ vì $167{;}167{;}168{;}168{;}168{;}168{;}169{;}169  \in [167;170)$
	 }
\end{ex}

% nguồn: Ok_TracNghiemLop11_Lan1
\begin{ex}%[1D5H1-1]%[ma:f4bfe170]
	Dựa trên đồ thị thể hiện thi đánh giá năng lực của  một trường đại học vào năm 2020 dưới đây.
	\begin{center}
		\begin{tikzpicture}[scale=1,>=stealth,font=\footnotesize,line join=round,line cap=round]
				\def\xmin{0} \def\xmax{13} 
				 \def\ymin{0} \def\ymax{9}
				 \draw[->] (\xmin,0)--(\xmax,0) node[below]{(Điểm)};
				 \draw[->] (0,\ymin)--(0,\ymax) node[above]{(Số học sinh)};
				\foreach  \y/\z in {1/20,2/40,3/60,4/80,5/100,6/120,7/140}{
					\draw  (0,\y) node[left]{$\z$};
				 }
				\foreach  \y/\z in {1/20,2/40,3/60,4/80,5/100,6/120,7/140}{
					\draw[gray!30](0,\y)--(13,\y); 
				} 
				 \foreach \x in {0.25,0.75,1.25,1.75,2.25,2.75,3.25,10.75,11.25,11.75}{
					 \draw (\x,0) node[above]{$0$};}
				\foreach \x/\t in {0.25/[0; 50],0.75/[51; 100],1.25/[101; 150],1.75/[151; 200],2.25/[201; 250],2.75/[251; 300],3.25/[301; 350],3.75/[351; 400],4.25/[401; 450],4.75/[451; 500],5.25/[501; 550],5.75/[551; 600],6.25/[601; 650],6.75/[651; 700],7.25/[701; 750],7.75/[751; 800],8.25/[801; 850],8.75/[851; 900],9.25/[901; 950],9.75/[951; 1000],10.25/[1001;  1050],10.75/[1051; 1100],11.25/[1101; 1150],11.75/[1151; 1200]}{
					\draw (\x,-0.775) node[rotate=90]{$\t$};} 
				\draw[fill=gray] (3.5,0)  rectangle (4,0.05)
					(3.75,0.05) node[above]{$1$}; 
				\draw[fill=gray] (4,0) rectangle  (4.5,0.05*8)
					(4.25,0.05*8) node[above]{$8$}; 
				\draw[fill=gray] (4.5,0) rectangle  (5,0.05*24)
					(4.75,0.05*24)  node[above]{$24$};
				\draw[fill=gray] (5,0) rectangle (5.5,0.05*54) 
					(5.25,0.05*54) node[above]{$54$}; 
				\draw[fill=gray] (5.5,0)  rectangle (6,0.05*95)
					 (5.75,0.05*95) node[above]{$95$};
				\draw[fill=gray]  (6,0) rectangle (6.5,0.05*95)
					 (6.25,0.05*95) node[above]{$95$};
				\draw[fill=gray] (6.5,0) rectangle (7,0.05*133) 
					 (6.75,0.05*133) node[above]{$133$};
				\draw[fill=gray] (7,0) rectangle  (7.5,0.05*122)
					(7.25,0.05*122)  node[above]{$122$};
				\draw[fill=gray] (7.5,0)  rectangle (8,0.05*104)
					(7.75,0.05*104) node[above]{$104$}; 
				\draw[fill=gray] (8,0) rectangle (8.5,0.05*62) 
					 (8.25,0.05*62) node[above]{$62$};
				\draw[fill=gray] (8.5,0)  rectangle (9,0.05*55)
					(8.75,0.05*55) node[above]{$55$}; 
				\draw[fill=gray] (9,0)  rectangle (9.5,0.05*21)
					(9.25,0.05*21)  node[above]{$21$};
				\draw[fill=gray] (9.5,0)  rectangle (10,0.05*12)
					(9.75,0.05*12) node[above]{$12$}; 
				\draw[fill=gray] (10,0) rectangle  (10.5,0.05*1)
					 (10.25,0.05*1) node[above]{$1$};
				\draw (\xmax/2,\ymax) node[above]{\textbf{Điểm thi đánh giá năng  lực của một trường đại học năm 2020}};
			\end{tikzpicture}
	\end{center}
	Tổng số học sinh tham  gia kì thi đánh giá năng lực trên là
	\choice 
		{$780$} 
		 {\True $787$}
		 {$696$}
		 {$697$}
	 \loigiai{
		Tổng số học sinh là  $1+8+24+54+95+95+133+122+104+62+55+21+12+1=787$ (học sinh).
	 }
\end{ex}

% nguồn: Ok_TracNghiemLop11_Lan1
\begin{ex}%[1D5H1-1]%[ma:becd5af5]
	Cho các số liệu thống kê ghi trong bảng sau đây. Thành tích chạy $50$ m của học sinh lớp $10A$ ở trường THPT B  (đơn vị: giây)
	\begin{center}
		\begin{tabular}{|r|r|r|r|r|r|r|}
			 \hline
			$6{,}3$ & $6{,}2$ & $6{,}5$  & $6{,}8$ & $6{,}9$ & $8{,}2$ & $8{,}6$ \\
			\hline 
			$6{,}6$ & $6{,}7$ & $7{,}0$ & $7{,}1$ & $7{,}2$ & $8{,}3$ &  $8{,}5$ \\
			 \hline
			$7{,}4$ & $7{,}3$ &  $7{,}2$ & $7{,}1$ & $7{,}0$ & $8{,}4$ & $8{,}1$ \\
			 \hline
			$7{,}1$  & $7{,}3$ & $7{,}5$ & $7{,}5$ & $7{,}6$ & $8{,}7$ & \\
			 \hline
			$7{,}6$  & $7{,}7$ & $7{,}8$ & $7{,}5$ & $7{,}7$ & $7{,}8$ & \\
			\hline 
		\end{tabular}
	\end{center}
	Trong lớp $10A$, số  học sinh chạy $50$ m hết $7$ giây đến dưới $8{,}5$ giây chiếm bao nhiêu phần trăm?
	 \choice
		{$30{,}30  \%$}
		{$27{,}27  \%$}
		 {$12{,}12 \%$}
		 {\True $69{,}69\%$}
	 \loigiai{
		Từ các số  liệu thống kê đã cho, ta xác định được:\\
		Tần số  của các lớp\\
		$n_{1}=2;  n_{2}=5; n_{3}=10; n_{4}=9; n_{5}=4; n_{6}=3$.\\
		Tần suất các lớp\\ 
		$f_{1}\approx 6{,}06 \%; f_{2}\approx 15{,}15 \%; f_{3}\approx 30{,}3 \%; f_{4}\approx 27{,}27 \%; f_{5}\approx  12{,}12 \%; f_{6}\approx 9{,}1 \%$.\\
		Từ đó ta  có bảng phân bố tần số ghép nhóm.\\
		Thành tích chạy $50$ m của học sinh lớp $10A$ ở trường  THPT B
		\begin{center}
			\begin{tabular}{|c|c|}
				\hline 
				Thời gian  chạy (giây) & Tần suất (\%) \\
				\hline 
				$[6{,}0; 6{,}5)$  & $6{,}06$ \\
				 \hline
				$[6{,}5; 7{,}0)$  & $15{,}15$ \\
				 \hline
				$[7{,}0; 7{,}5)$ & $30{,}30$  \\
				 \hline
				 $[7{,}5; 8{,}0)$ & $27{,}27$ \\
				 \hline
				$[8{,}0; 8{,}5)$ & $12{,}12$  \\
				\hline 
				 $[8{,}5; 9{,}0)$ & $9{,}10$ \\
				 \hline
				Cộng &  $100(\%)$ \\
				\hline 
			\end{tabular}
		\end{center}
	Suy ra số học sinh chạy $50$ m  hết $7$ giây đến dưới $8{,}5$ giây chiếm: \[30{,}30 \%+27{,}27 \%+12{,}12 \%=69{,}69 \%.\]}
\end{ex}

% nguồn: Ok_TracNghiemLop11_Lan1
\begin{ex}%[1D5H1-1]%[ma:245bd174]
	Nhóm số liệu rời rạc $k_1-k_2$ với $k_1$,  $k_2\in \mathbb{N}$, $k_1<k_2$ là nhóm gồm các giá trị
	 \choice
		 {$k_1$ và $k_2$}
		{$k+1$, $\ldots$,  $k_2$}
		{$k_1$, $\ldots$,  $k_2+1$}
		 {\True $k_1$, $k_1+1$, $\ldots$, $k_2$}
	 \loigiai{
		Với $k_1$, $k_2\in \mathbb{N}$ và $k_1<k_2$ suy ra nhóm số liệu rời rạc từ $k_1$ đến $k_2$ gồm $k_1$, $k_1+1$, $k_1+2$, $\ldots$,  $k_2$.
	} 
\end{ex}

% nguồn: Ok_TracNghiemLop11_Lan1
\begin{ex}%[1D5H1-1]%[ma:ccc79f16]
	Khảo sát thời gian tập thể dục trong ngày của một số học sinh khối 11 thu được mẫu số liệu ghép nhóm sau 
	\begin{center}
		\begin{tabular}{|c|c|c|c|c|c|}
			\hline 
			 Thời gian (phút)&$[0;20)$&$[20;40)$&$[40;60)$&$[60;80)$&$[80;100)$\\
			\hline 
			Số  học sinh &$5$&$9$&$12$&$10$&$6$\\
			\hline 
		\end{tabular}
	\end{center}
	Giá trị đại diện của  nhóm $[20; 40)$ là
	 \choice
		{$10$} 
		 {$20$}
		{\True  $30$}
		{$40$} 
	 \loigiai{
		Trong mỗi khoảng  thời gian, giá trị đại diện là trung bình cộng của giá trị hai đầu mút nên ta có giá trị đại diện của nhóm $[20;40)$ là $\dfrac{20+40}{2}=30$.
	 }
\end{ex}

% nguồn: Ok_TracNghiemLop11_Lan1
\begin{ex}%[1D5H1-1]%[ma:4ed72832]
	Giá trị đại diện của nhóm $[a_i;a_{i+1})$  là
	 \choice
		{$a_i$} 
		{$a_{i+1}$} 
		{$\dfrac{a_{i+1}-a_i}{2}$} 
		{\True $\dfrac{a_{i+1}+a_i}{2}$} 
	 \loigiai{
		Giá trị đại diện của  nhóm $[a_i;a_{i+1})$ là $\dfrac{a_{i+1}+a_i}{2}$.
	 }
\end{ex}

% nguồn: Ok_TracNghiemLop11_Lan1
\begin{ex}%[1D5B1-1]%[ma:73ec34c1]
	Tính số gia $\Delta y$ của  hàm số $y = x^3 - 2x$ theo số gia của đối số $\Delta x$ tại $x = 1$.
	 \choice
		{$\Delta y  = (\Delta x)^3 + 3 (\Delta x)^2 + 3 \Delta x + 1$}
		{$\Delta  y = (\Delta x)^3 - 3 (\Delta x)^2 + \Delta x$}
		{\True $\Delta y = (\Delta x)^3 +  3 (\Delta x)^2 + \Delta x$}
		{$\Delta y = (\Delta x)^3 - 3 (\Delta x)^2 + 3 \Delta x +  1$}
	\loigiai{ 
	 Ta có
	\begin{align*}
		\Delta y & = \left(1 + \Delta  x \right)^3 - 2(1 + \Delta x) - (1^3 - 2 \cdot 1) \\
		& = 1 + 3 \Delta x + 3 (\Delta x)^2 + (\Delta x)^3 - 2 - 2\Delta x + 1 \\ 
		& = (\Delta  x)^3 + 3 (\Delta x)^2 + \Delta x.
	\end{align*}
	 }
\end{ex}

% nguồn: Ok_TracNghiemLop11_Lan1
\begin{ex}%[1D5Y1-1]%[ma:57f6bdf5]
	Trong các phát biểu sau, phát biểu nào sau đây là  \textbf{đúng}?
	 \choice
		{Nếu hàm số $y=f(x)$ không  liên tục tại $x_0$ thì nó có đạo hàm tại điểm đó}
		{Nếu hàm số $y=f(x)$ có đạo hàm tại $x_0$ thì nó không liên tục tại  điểm đó}
		{\True Nếu hàm số $y=f(x)$ có đạo hàm tại $x_0$  thì nó liên tục tại điểm đó}
		{Nếu hàm số $y=f(x)$ liên tục tại $x_0$ thì nó có  đạo hàm tại điểm đó}
	 \loigiai{
		\lq\lq Nếu hàm số  $y=f(x)$ có đạo hàm tại $x_0$ thì nó liên tục tại điểm đó \rq\rq\, là mệnh đề \textbf{đúng}.\\
		Điều ngược lại chưa chắc đúng tức là \lq\lq Nếu hàm số $y=f(x)$ liên tục tại $x_0$ thì nó có đạo hàm tại điểm đó \rq\rq\, là mệnh đề  là mệnh đề \textbf{sai}.\\
		\textbf{Ví  dụ:} Hàm số $y=|x|$ liên tục tại $x=0$ nhưng không có đạo hàm tại $x=0$.
	} 
\end{ex}

% nguồn: Ok_TracNghiemLop11_Lan1
\begin{ex}%[1D5Y1-1]%[ma:b969e92e]
	Tính đạo hàm của hàm số $y=k$, với $k$ là  hằng số.
	\choice 
		{$y'=-1 $} 
		{\True $y'=0 $} 
		{$y'=1  $}
		{$y'=k  $}
	\loigiai{Ta có $y'=0  $.
	} 
\end{ex}

% nguồn: Ok_TracNghiemLop11_Lan1
\begin{ex}%[1D5B1-1]%[ma:f3bcf798]
	Tiếp tuyến với đồ thị của hàm số $f(x) = \dfrac{3}{{2x - 1}}$ tại điểm  có hoành độ $x = 2$ có hệ số góc là \\
	\choice{\True $ - \dfrac{2}{3}$}{$\dfrac{2}{3}$}{$ 2  $}{$ -2 $}
	\loigiai{Tính $  f'(2)=\dfrac{-2}{3} $.
	} 
\end{ex}

% nguồn: Ok_TracNghiemLop11_Lan1
\begin{ex}%[1D5Y1-1]%[ma:5868a78d]
	Cho hàm số $ f\left( x \right)  $ liên tục trên khoảng $ (a;b) $, $ x_0 \in (a;b) $. Tính $ f'\left( x_0 \right) $ bằng định nghĩa ta cần tính
	 \choice
		 {$ \mathop {\lim }\limits_{x \to 0} \dfrac{{\Delta y}}{{\Delta x}} $}
		 {\True $ \mathop {\lim }\limits_{\Delta x \to 0} \dfrac{{\Delta y}}{{\Delta x}} $}
		{$ \mathop {\lim }\limits_{\Delta x \to 0}  \dfrac{{\Delta x}}{{\Delta y}} $}
		{$ \mathop {\lim }\limits_{x \to 0}  \dfrac{{\Delta y}}{{x}} $}
	\loigiai{ 
		Theo định nghĩa đạo hàm ta có  $ f'\left( x_0 \right) = \mathop {\lim }\limits_{\Delta x \to 0} \dfrac{{\Delta y}}{{\Delta x}} $.
	} 
\end{ex}

% nguồn: Ok_TracNghiemLop11_Lan1
\begin{ex}%[1D5B1-1]%[ma:10bd6a20]
	Tính  tỷ số $\dfrac{\Delta y}{\Delta x}$ của hàm số $y=f(x)=2x^3$ theo $x$ và $\Delta x$.
	\choice 
		{$\dfrac{\Delta y}{\Delta x}=\dfrac{2x^3-2{{\left( \Delta x \right)}^3}}{\Delta  x}$}
		{$\dfrac{\Delta  y}{\Delta x}=2{{\left( \Delta x \right)}^2}$}
		{\True $\dfrac{\Delta y}{\Delta x}=6x^2+6x\Delta x+2{{\left( \Delta x \right)}^2}$} 
		{$\dfrac{\Delta  y}{\Delta x}=3x^2+3x\Delta x+{{\left( \Delta x \right)}^2}$}
	 \loigiai{
		Ta có $$\Delta  y=f\left( x+\Delta x \right)-f(x)=2{{\left( x+\Delta x \right)}^3}-2x^3=6x^2\Delta x+6x{{\left( \Delta x \right)}^2}+2{{\left( \Delta x \right)}^3}$$
		Suy ra $$\dfrac{\Delta y}{\Delta  x}=6x^2+6x\Delta x+2{{\left( \Delta x \right)}^2}.$$
	 }
\end{ex}

% nguồn: Ok_TracNghiemLop11_Lan1
\begin{ex}%[1D5B1-1]%[ma:5dfc6e17]
	Hàm số $f(x_0)$ có đạo hàm tại $x_0$  là $f'(x_0)$. Chọn khẳng định \textbf{sai}.
	 \choice
		{$f'(x_0) = \lim\limits_{x\to x_0} \dfrac{f(x)  - f(x_0)}{x-x_0}$}
		{$f'(x_0)  = \lim\limits_{\Delta x\to 0} \dfrac{f(x_0 + \Delta x) - f(x_0)}{\Delta x}$}
		{$f'(x_0) = \lim\limits_{h\to 0} \dfrac{f(x_0 + h)  - f(x_0)}{h}$}
		{\True $f'(x_0) = \lim\limits_{x\to x_0} \dfrac{f(x  + x_0)-f(x_0)}{x-x_0}$}
	\loigiai{ 
		\begin{itemize}
			\item Theo định nghĩa đạo hàm tại một điểm SGK ĐS-GT 11,  ta có
			$f'(x_0) = \lim\limits_{x\to x_0}  \dfrac{f(x) - f(x_0)}{x-x_0}.$
			\item Đặt $\Delta  x = x - x_0$, ta có
			$$f'(x_0) = \lim\limits_{x\to x_0} \dfrac{f(x) - f(x_0)}{x-x_0} = \lim\limits_{\Delta x\to 0} \dfrac{f(x_0 + \Delta x) - f(x_0)}{x_0 + \Delta x - x_0} = \lim\limits_{x\to  x_0} \dfrac{f(x) - f(x_0)}{x-x_0}.$$
			\item Đặt $h = x - x_0$, ta có 
			$$f'(x_0) = \lim\limits_{x\to x_0} \dfrac{f(x) - f(x_0)}{x-x_0} = \lim\limits_{h \to 0} \dfrac{f(x_0 + h)  - f(x_0)}{x_0 + h - x_0} = \lim\limits_{h\to 0} \dfrac{f(x_0 + h) - f(x_0)}{h}.$$
			\item Như vậy, đáp án sai là $f'(x_0) = \lim\limits_{x\to x_0} \dfrac{f(x +  x_0)-f(x_0)}{x-x_0}$.
		\end{itemize}
	 }
\end{ex}

% nguồn: Ok_TracNghiemLop11_Lan1
\begin{ex}%[1D5Y1-1]%[ma:a43ea2ac]
	 Khẳng định nào sau đây là đúng?
	\choice 
		{\True $ f'(x_0)=\displaystyle  \lim_{x\to x_0}{\dfrac{f(x)-f(x_0)}{x-x_0}}
		 $}
		{$ f'(x_0)=\displaystyle  \lim_{x\to x_0}{\dfrac{f(x)+f(x_0)}{x-x_0}}
		$  }
		{$  f'(x_0)=\displaystyle \lim_{x\to x_0}{\dfrac{f(x_0)-f(x)}{x-x_0}}
		 $ }
		{$ f'(x_0)=\displaystyle \lim_{x\to  0}{\dfrac{f(x)-f(x_0)}{x-x_0}}
		$  }
	 \loigiai{
	 }
\end{ex}

% nguồn: Ok_TracNghiemLop11_Lan1
\begin{ex}%[1D5Y1-1]%[ma:88290149]
	Cho hàm số $y=f(x)$ có đạo hàm thỏa  mãn $f'(6)=2$. Tính giá trị của biểu thức $\lim\limits_{ x \to 6}\dfrac{ f(x)-f(6)}{x-6}$.
	 \choice
		{\True $2$} 
		{$\dfrac{1}{3}$} 
		 {$\dfrac{1}{2}$}
		{$12$} 
	\loigiai{ 
	Theo định nghĩa đạo hàm của hàm số tại $x=6$, suy ra  $\lim\limits_{ x \to 6}\dfrac{ f(x)-f(6)}{x-6} =f'(6)=2$.}
\end{ex}

% nguồn: Ok_TracNghiemLop11_Lan1
\begin{ex}%[1D5B1-1]%[ma:afb688ee]
	Cho hàm số  $f(x)$ xác định bởi $f(x)=\heva{\dfrac{\sqrt{x^2+1}-1}{x} &(x\ne 0)\\ 0\quad &(x=0)}$. Giá trị $f'(0)$ là
	 \choice
		 {0}
		{Không  tồn tại}
		{\True $\dfrac{1}{2}$} 
		 {1}
	 \loigiai{
		$f'(0)=\lim\limits_{x\to 0}\dfrac{f(x)-f(0)}{x-0}=\lim\limits_{x\to  0}\dfrac{\sqrt{x^2+1}-1}{x^2}\overset{\text{nhân liên hợp}}{=}\lim\limits_{x\to 0}\dfrac{1}{\sqrt{x^2+1}+1}=\dfrac{1}{2}$.
	 }
\end{ex}

% nguồn: Ok_TracNghiemLop11_Lan1
\begin{ex}%[1D5B1-1]%[ma:585d1993]
	Cho hàm số $ f\left( x  \right) $ có đạo hàm trên tập số thực. Tìm hệ thức đúng.
	 \choice
		{\True $ f'\left( {1} \right) = \mathop {\lim }\limits_{x \to 1}  \dfrac{{f\left( x \right) - f\left( 1 \right)}}{{x - 1}}$}
		 {$ f'\left( {1} \right) = \mathop {\lim }\limits_{x \to 1} \dfrac{{f\left( x \right) }}{{x - 1}}$}
		{$ f'\left( {1} \right) = \mathop {\lim }\limits_{x \to 1} \dfrac{{f\left(  x \right)}}{{x}}$}
		{$ f'\left( {1} \right) = \mathop {\lim }\limits_{x  \to 1} \dfrac{{f\left( 1 \right)}}{{x - 1}}$}
	\loigiai{ 
		Ta  có $ f'\left( {x_0} \right) = \mathop {\lim }\limits_{x \to x_0} \dfrac{{f\left( x \right) - f\left( x_0 \right)}}{{x - x_0}}$.
	 }
\end{ex}

% nguồn: Ok_TracNghiemLop11_Lan1
\begin{ex}%[1D5B1-1]%[ma:662d1060]
	Có bao nhiêu mệnh đề đúng trong  các mệnh đề sau?
	\begin{enumerate}[(1)]
		\item Nếu hàm số $f(x)$  có đạo hàm tại điểm $x = x_0$ thì $f(x)$ liên tục tại $x_0$.
		\item Nếu hàm số $f(x)$ liên tục tại điểm $x = x_0$ thì $f(x)$ có  đạo hàm tại $x_0$.
		\item  Nếu hàm số $f(x)$ gián đoạn tại điểm $x = x_0$ thì $f(x)$ không có đạo hàm tại $x_0$.
	\end{enumerate}
	\choice 
		{$0$} 
		 {$1$}
		{\True  $2$}
		 {$3$}
	\loigiai{ 
		Hàm số liên tục tại một điểm có thể không có đạo hàm tại điểm  đó. Như vậy, chỉ có hai mệnh đề đúng là mệnh đề $(1)$ và $(3)$.
	} 
\end{ex}

% nguồn: Ok_TracNghiemLop11_Lan1
\begin{ex}%[1D5Y1-1]%[ma:036f8fe4]
	Cho hàm số $f(x)$ có đạo hàm  tại $x_0$. Chọn khẳng định \textbf{đúng} trong các khẳng định dưới đây?
	\choice 
		 {$f(x)$ gián đoạn tại $x_0$}
		{$f(x)$ không  cần xác định tại $x_0$}
		{$\lim\limits_{x \to x_0}  f(x)=0$}
		{\True $f(x)$ liên tục tại $x_0$} 
	 \loigiai{
		$f(x)$ gián đoạn tại $x_0$ sai vì theo Định lý 1 trang 150 SGK Đại Số và Giải Tích 11 cơ bản, nếu hàm số $f(x)$ có đạo hàm tại $x_0$ thì nó liên tục tại điểm đó. Như vậy $f(x)$  liên tục tại $x_0$ đúng.\\
		$f(x)$ không cần xác định tại $x_0$ sai vì theo định nghĩa đạo hàm của hàm số tại một điểm SGK  Đại Số và Giải Tích 11 cơ bản, trang 148, hàm số $f(x)$ phải xác định tại $x_0$.\\
	$\lim\limits_{x \to x_0} f(x)=0$ sai. Chẳng hạn, xét hàm số $f(x)=x^2$. Đạo hàm của $f(x)=x^2$ tại  $x_0=1$ tồn tại và bằng $2$. Nhưng $\lim\limits_{x \to 1} x^2=1 \ne 0$.}
\end{ex}

% nguồn: Ok_TracNghiemLop11_Lan1
\begin{ex}%[1D5B1-1]%[ma:941a6412]
	Cho hàm số $y=\dfrac{5}{3}{x^3}-x^2+4$ có đồ thị  $(C)$. Tiếp tuyến của $(C)$ tại điểm có hoành độ $x_0=3$ có hệ số góc là
	\choice 
		 {\True $39$}
		 {$40$}
		 {$51$}
		 {$3$}
	\loigiai{ 
		Hệ số góc của tiếp  tuyến của đồ thị hàm số $y=f(x)$ tại điểm có hoành độ $x_0$ là $k=f'(x_0)$.\\
		Hàm số $y=\dfrac{5}{3}{x^3}-x^2+4$ có  $y'=5x^2-2x$.\\
	Suy ra tiếp tuyến của $(C)$ tại điểm có hoành độ $x_0=3$ có hệ số góc $k=y'(3)=39$.} 
\end{ex}

% nguồn: Ok_TracNghiemLop11_Lan1
\begin{ex}%[1D5Y1-1]%[ma:8ee12e04]
	Cho hàm số $y=f(x)$ có  đạo hàm tại $x_0$ là $f'(x_0)$. Khẳng định nào sau đây {\bf sai}?
	\choice 
		{$f'(x_0)=\lim\limits_{x \to x_0}\dfrac{f(x)-f(x_0)}{x-x_0}$} 
		{\True $f'(x_0)=\lim\limits_{x  \to x_0}\dfrac{f\left(x+x_0\right)-f(x_0)}{x-x_0}$}
		{$f'(x_0)=\lim\limits_{h \to  0}\dfrac{f\left(x_0+h\right)-f(x_0)}{h}$}
		 {$f'(x_0)=\lim\limits_{\Delta x\to 0}\dfrac{f\left(x_0+\Delta x\right)-f(x_0)}{\Delta x}$}
	\loigiai{ 
		Định nghĩa: Cho hàm số $y=f(x)$ xác định trên $\left(a;b\right)$  và $x_0\in (a;b)$. Giới hạn hữu hạn (nếu có) của tỉ số $\dfrac{f(x)-f(x_0)}{x-x_0}$ khi $x$ dần đến $x_0$ gọi là đạo hàm của hàm số đã cho tại điểm $x_0$, kí hiệu là $f'(x_0)$, ta có $f'(x_0)=\lim\limits_{x \to x_0} \dfrac{f(x)-f(x_0)}{x-x_0}$.\\
		Từ định nghĩa rút ra kết luận đáp  án \lq\lq$f'(x_0)=\lim\limits_{x \to x_0} \dfrac{f(x+x_0)-f(x_0)}{x-x_0}$\rq\rq\, sai.\\
		\lq\lq $f'(x_0)=f'(x_0)=\lim\limits_{x \to x_0}\dfrac{f(x)-f(x_0)}{x-x_0}$\rq\rq\,  đúng do định nghĩa.\\
		\lq\lq  $f'(x_0)=\lim\limits_{h \to 0}\dfrac{f(x_0+h)-f(x_0)}{h}$\rq\rq\, đúng vì đặt $x=x_0+h\Rightarrow \heva{& x-x_0=h \\ & x\to x_0\Rightarrow h\to 0}$.\\
		\lq\lq $f'(x_0)=\lim\limits_{\Delta x\to 0}\dfrac{f(x_0+\Delta x)-f(x_0)}{\Delta x}$\rq\rq\, đúng vì đặt $x=x_0+\Delta x\Rightarrow \heva{& x-x_0=\Delta x \\ &  x\to x_0\Rightarrow \Delta x\to 0}$.
	 }
\end{ex}

% nguồn: Ok_TracNghiemLop11_Lan1
\begin{ex}%[1D5Y1-1]%[ma:410ee450]
	Tính số gia của hàm số $y=f(x)=x^3+x^2+1$ tại điểm $x_0$ ứng với  số gia $\Delta x=1$.
	\choice 
		{$\Delta y=3x_0^2+5x_0+3$} 
		{$\Delta y=2x_0^3+3x_0^2+5x_0+2$} 
		{\True $\Delta y=3x_0^2+5x_0+2$} 
		{$\Delta y=3x_0^2-5x_0+2$} 
	\loigiai{ 
		Ta có 
		\begin{eqnarray*}
			\Delta y &=&f\left( x_0+\Delta x \right)-f\left(  x_0 \right)=f\left( x_0+1 \right)-f\left( x_0 \right)\\
			&=&\left[\left( x_0+1 \right)^3+\left(x_0+1\right)^2+1 \right]-\left[x_0^3+x_0^2+1 \right]\\ 
			 &=&3x_0^2+5x_0+2.
		\end{eqnarray*}
	} 
\end{ex}

% nguồn: Ok_TracNghiemLop11_Lan1
\begin{ex}%[1D5B1-1]%[ma:27a2804f]
	Cho hàm số $y=f(x)$ xác định trên $\mathbb{R}$ thỏa mãn $ \lim\limits_{x \to 2} \dfrac{f(x)-f(2)}{x-2}=3$. Kết  quả nào sau đây là đúng?
	\choice 
		 {$f'(3)=2$}
		{\True  $f'(2)=3$}
		{$f'(x)=3$} 
		{$f'(x)=2$} 
	\loigiai{Theo định nghĩa đạo hàm ta có  $f'(x_0)= \lim\limits_{x \to x_0} \dfrac{f(x)-f(x_0)}{x-x_0}$.\\
	Dó đó $$ \lim\limits_{x  \to 2} \dfrac{f(x)-f(2)}{x-2}=f'(2)=3.$$}
\end{ex}

% nguồn: Ok_TracNghiemLop11_Lan1
\begin{ex}%[1D5B1-1]%[ma:9c0d5c24]
	Cho hàm số $y=f(x)$ xác định trên  thỏa mãn $\lim\limits_{x\to 3}\dfrac{f(x)-f(3)}{x-3}=2$. Kết quả đúng là
	\choice 
		 {$f'(2)=3$}
		 {$f'(x)=2$}
		 {$f'(x)=3$}
		{\True $f'(3)=2$} 
	 \loigiai{
		Theo định nghĩa đạo hàm của hàm  số tại một điểm ta có\\
	 \[\lim\limits_{x\to 3}\dfrac{f(x)-f(3)}{x-3}=2=f'(3).\]}
\end{ex}

% nguồn: Ok_TracNghiemLop11_Lan1
\begin{ex}%[1D5B1-1]%[ma:097f39e5]
	Tìm số gia $ \Delta y $ của hàm số $ y = x^2 $ biết $ x_0 =  3 $ và $ \Delta x = -1 $.
	 \choice
		{$  \Delta y = 13 $}
		 {$ \Delta y = 7 $}
		{\True $ \Delta y  = -5 $}
		{$ \Delta y = 16 $} 
	 \loigiai{
		Ta có $ \Delta y  = y\left( {x_0 + \Delta x} \right) - y\left( {x_0 } \right) = y\left( 2 \right) - y\left( 3 \right) = 2^2 - 3^2 = - 5$.
	 }
\end{ex}

% nguồn: Ok_TracNghiemLop11_Lan1
\begin{ex}%[1D5B1-1]%[ma:bc7b7e99]
	Cho các khẳng định sau, có bao nhiêu khẳng định \textbf{đúng}? 
	\begin{enumerate}
		\item Một hàm số liên tục tại  một điểm có thể không có đạo hàm tại điểm đó.
		\item Nếu hàm số liên tục tại $ x_0  $ thì nó có đạo hàm tại điểm đó.
		\item Nếu hàm số  $ y= f(x) $ có đạo hàm tại $ x_0 $ thì nó liên tục tại điểm đó.
		\item Nếu hàm số $ y = f(x) $ gián  đoạn tại $ x_0 $ thì nó không có đạo hàm tại điểm đó.
	\end{enumerate}
	 \choice
		{$ 1  $}
		 {$ 2 $}
		{\True $ 3 $} 
		 {$ 4 $}
	\loigiai{ 
		Theo định lý về quan hệ giữa sự tồn tại của đạo hàm và tính liên tục của hàm số,  khẳng định a), c), d) đúng.
	} 
\end{ex}

% nguồn: Ok_TracNghiemLop11_Lan1
\begin{ex}%[1D5B1-1]%[ma:03823bb5]
	Số  gia $\Delta y$ của hàm số $y = x^2 + 2x - 5$ tại điểm $x_0 = 1$ là
	\choice 
		{${\left( {\Delta x} \right)^2} - 4\Delta  x$}
		{\True ${\left( {\Delta x} \right)^2} + 4\Delta  x$}
		{${\left( {\Delta  x} \right)^2} - 2\Delta x$}
		{${\left(  {\Delta x} \right)^2} + 2\Delta x - 5$}
	\loigiai{$ \Delta y =f(1 + \Delta x) - f(1)=(1 + \Delta x)^2 + 2(1 + \Delta x) - 5  - (-2)={\left( {\Delta x} \right)^2} + 4\Delta x$.
	} 
\end{ex}


%%%%%%%%%%%%%%%% ĐÚNG/SAI %%%%%%%%%%%%%%%%

% nguồn: Ok_dungsai_lop11_Lan2
\begin{ex}%[1D5V1-1]%[ma:3982cca2]
	 Điểm kiểm tra giữa học kì $ 1 $ của lớp 11A1 được thống kê như bảng dưới
	\begin{center}
		\begin{tabular}{|c|c|c|c|c|c|}
			\hline 
			Điểm & $[6;7)$ & $[6;7)$ & $[7;8)$ & $[8;9)$ & $[9;10)$  \\
			 \hline
			Số học sinh & $ 3 $ & $ 5 $ & $ 10 $ & $6 $ & $ 2 $  \\
			\hline 
		\end{tabular}
	\end{center}
	 \choiceTF
		{Cỡ  mẫu của số liệu là $n=25$}
		{\True Nhóm  chứa mốt của số liệu là $[7;8)$}
		{Số học sinh đạt điểm  khá (từ $ 8 $ điểm trở lên) là $ 18 $}
		{\True Tứ phân vị thứ nhất của mẫu số liệu trên là:  $Q_1=6{,}7$}
	 \loigiai{
		\begin{itemchoice}
			\itemch \textbf{Sai.} \\ 
			 Cỡ mẫu số liệu là $n=3+5+10+6+2=26$.
			 \itemch \textbf{Đúng.} \\
			Nhóm có  tần số lớn nhất là $[7;8)$.
			 \itemch \textbf{Sai.} \\
			Số học sinh đạt điểm khá (từ $  8 $ điểm trở lên) là $6+2=8$.
			 \itemch \textbf{Đúng. }\\
			 Cỡ mẫu: $n=26$.\\
			Gọi $x_1, x_2,\cdots, x_{26}$ là mẫu số liệu được sắp xếp theo  thứ tự không giảm.\\
			Trung  vị của mẫu số liệu: $\dfrac{x_{13}+x_{14}}{2} \in [7;8)$.\\
			Xét  nửa mẫu số liệu bên trái $x_1, x_2,cldots, x_{13}$ có trung vị $x_7 \in [6;7)$.\\
			 Ta có: $p=2$ , $a_2=6$ , $a_3=7$ , $m_2=5$ , $m_1=3$ , $n=26$. Khi đó $ Q_1 = 6{,}7 $.
		\end{itemchoice}
	 }
\end{ex}

% nguồn: Ok_dungsai_lop11_Lan2
\begin{ex}%[1D5V1-1]%[ma:f34d15fc]
	Thời gian truy cập Internet mỗi buổi tối của một số học sinh được cho trong bảng  sau:
	\begin{center}
		\begin{tabular}{|c|c|c|c|c|c|}
			 \hline
			Thời gian (phút) & $ [9{,}5;12{,}5) $ &$ [12{,}5;15{,}5) $ & $ [15{,}5;18{,}5) $ & $ [18{,}5;21{,}5) $ & $ [21{,}5;24{,}5) $  \\
			\hline 
			 Số học sinh & $ 3 $ & $ 12 $ & $ 15 $ & $ 24 $ & $ 2 $ \\
			 \hline
		\end{tabular}
	\end{center}
	 \choiceTF
		{\True Số học sinh truy cập Internet mỗi buổi tối có thời gian  từ $ 18{,}5 $ phút đến dưới $ 21{,}5 $ phút là $24$}
		{\True Nhóm chứa mốt là $[18{,}5;21{,}5)$} 
		{Số trung bình của mẫu số liệu trên thuộc  khoảng $[18{,}5;21{,}5)$}
		{Mốt của mẫu  số liệu trên thuộc khoảng $[18{,}5;21{,}5)$}
	 \loigiai{
		\begin{itemchoice}
			\itemch \textbf{Đúng.}\\ 
			Số học sinh truy cập Internet mỗi buổi tối có thời gian từ $ 18{,}5 $ phút đến dưới $ 21{,}5 $  phút là $24$.
			 \itemch \textbf{Đúng.}\\
			Nhóm chứa  mốt là $[18{,}5;21{,}5)$.
			 \itemch \textbf{Sai.}\\
			Số trung bình của mẫu số  liệu bằng $$(11\cdot 3+14\cdot 12+17\cdot 15+20\cdot 24+23\cdot 2):56=\dfrac{491}{28} \approx 17{,}5.$$
			\itemch  \textbf{Sai.}
			\begin{center}
				\begin{tabular}{|c|c|c|c|c|c|}
					 \hline
					Thời gian (phút) & $ [9{,}5;12{,}5)  $ &$ [12{,}5;15{,}5) $ & $ [15{,}5;18{,}5) $ & $ [18{,}5;21{,}5) $ & $ [21{,}5;24{,}5) $ \\
					 \hline
					Giá  trị đại diện & $ 11 $ & $ 14 $ & $ 17 $ & $ 20 $ & $ 23 $ \\
					 \hline
					Số học sinh  & $ 3 $ & $ 12 $ & $ 15 $ & $ 24 $ & $ 2 $ \\
					 \hline
				\end{tabular}
			\end{center}
			Nhóm chứa mốt là $[18{,}5;21{,}5)$.\\ 
			Ta có  $u_m=18{,}5$; $n_{m-1}=15$; $n_m=24$; $n_{m+1}=2$; $u_{m+1}-u_m=3$.\\
			 Mốt của mẫu số liệu ghép nhóm là\\ $M_o=18+\dfrac{24-15}{(24-15)+(24-2)}\cdot 3=\dfrac{585}{31} \approx 18{,}8$.
		\end{itemchoice}
	 }
\end{ex}

% nguồn: Ok_dungsai_lop11_Lan2
\begin{ex}%[1D5V1-1]%[ma:843b1878]
	Cân nặng  (kg) của một số học sinh lớp $11$ được cho như bảng sau
	\begin{center}
		\begin{tabular}{|c|c|c|c|c|c|}
			 \hline
			Cân nặng &  $[45;49)$ &$[49;53)$ & $[53;57)$ & $[57;61)$ & $[61;65)$ \\
			\hline 
			Số học sinh & $5$ & $7$ & $11$ & $9$ &  $6$ \\
			\hline 
		\end{tabular}
	\end{center}
	 \choiceTF
		{\True  Cỡ mẫu $n=38$}
		{Giá trị đại diện của nhóm $[49;53)$  là $52$}
		{\True Nhóm chứa mốt của mẫu số  liệu là $[53;57)$}
		{ Cân nặng trung  bình của mỗi học sinh là $55$}
	\loigiai{ 
		\begin{itemchoice}
			\itemch  \textbf{Đúng.}\\
			Cỡ mẫu $n=5+7+11+9+6=38$. 
			 \itemch \textbf{Sai.}\\
			Giá  trị đại diện của nhóm $[49;53)$ bằng $\dfrac{1}{2} (49+53)=51$.
			\itemch \textbf{Đúng.}\\ 
			Nhóm $[53;57)$ có tần số  $11$.
			\itemch \textbf{Sai.} \begin{center}
			\begin{tabular}{|c|c|c|c|c|c|}
				\hline 
				Cân nặng  & $ [45;49) $ &$ [49;53) $ & $ [53;57) $ & $ [57;61) $ & $ [61;65) $ \\
				 \hline
				Giá trị  đại diện & $ 47 $ & $ 51 $ & $ 55 $
				& $ 59 $ & $  63 $\\
				 \hline
				Số học sinh & $ 5 $ & $ 7  $ & $ 11 $ & $ 9 $ & $ 6 $ \\
				 \hline
			\end{tabular}
		\end{center}
		 Cân nặng trung bình của mỗi học sinh bằng\\ $(47\cdot 5+51\cdot 7+55\cdot 11+59\cdot 9+63\cdot 6):38\approx 55{,}42$.
	\end{itemchoice}
} 
\end{ex}

% nguồn: Ok_dungsai_lop11_Lan2
\begin{ex}%[1D5V1-1]%[ma:c5738983]
	Một nhà thực vật học đô chiều dài trung bình  của $ 74 $ lá cây (đơn vị: milimét) và thu được bảng tần số ghép nhóm như sau
	\begin{center}
		\begin{tabular}{|c|c|c|}
			\hline 
			Nhóm &  Giá trị đại diện & Tần số \\
			 \hline
			$[5{,}45;5{,}85)$ & $  5{,}65 $ & $ 5 $ \\
			 \hline
			$[5{,}85;6{,}25)$ &  $ 6{,}05 $ & $ 9 $ \\
			\hline 
			$[6{,}25;6{,}65)$ & $ 6{,}45 $ & $ 15 $ \\ 
			 \hline
			$[6{,}65;7{,}05)$ &  $ 6{,}85 $ & $ 19 $ \\
			 \hline
			$[7{,}05;7{,}45)$ & $ 7{,}25 $ & $ 16 $  \\
			\hline 
			$[7{,}45;7{,}85)$  & $ 7{,}65 $ & $ 8 $ \\
			 \hline
			$[7{,}85;8{,}25)$ & $ 8{,}05 $ & $ 2 $  \\
			\hline 
		\end{tabular}
	\end{center}
	 \choiceTF
		{Chiều dài trung bình của $ 74 $ lá cây gần bằng  $ 6{,}4$ (\; mm)}
		 {\True Độ dài nhóm là $0{,}4$}
		{Nhóm chứa mốt là: $[7{,}05;7{,}45)$} 
		{Mốt của  mẫu số liệu ghép nhóm là $ 6{,}65$}
		\begin{flushright}
			\textit{(Các kết quả tính  được trong bài làm tròn đến hàng phần trăm)}
	\end{flushright}
	 \loigiai{
		\begin{itemchoice}
			 \itemch \textbf{Sai.}\\
			Chiều dài trung bình của $ 74 $  lá cây là\\
			\[\overline{x} =\dfrac{5{,}65.5+6{,}05.9+6{,}45\cdot 15+6{,}85\cdot 19+7{,}25.16+7{,}65\cdot  8+8{,}05.2}{74} =\dfrac{5029}{740} \approx 6{,}8\text{(mm).}
			\] 
			\itemch  \textbf{Đúng.} \\
			Độ dài nhóm là  $0{,}4$.
			 \itemch \textbf{Sai.} \\
			Nhóm chứa  mốt là $[6{,}65;7{,}05)$.
			\itemch  \textbf{Sai.} \\
			Do nhóm chứa mốt $[6{,}65;7{,}05)$,  trong đó
			$u_m=6{,}65$; $u_{m+1}=7{,}05$; $u_{m+1}-u_m=0{,}4$;  $n_m=19$; $n_{m-1}=15$; $n_{m+1}=16.
			 $\\
			 Do đó mốt của mẫu số liệu ghép nhóm là\\
			\[M_o=6{,}65+\dfrac{19-15}{(19-15)+(19-16)} \cdot 0{,}4\approx  6{,}88.
			 \]
		\end{itemchoice}
	 }
\end{ex}

% nguồn: Ok_dungsai_lop11_Lan2
\begin{ex}%[1D5V1-1]%[ma:9c8579c9]
	Thảo thống kê lại số bước chân bạn đi mỗi  ngày trong $ 3 $ tháng. Kết quả được tổng hợp ở bảng sau
	\begin{center}
		\begin{tabular}{|c|c|}
			\hline 
			Số bước chân & Số  ngày \\
			\hline 
			$[4\,000;  5\,000)$ & $12$ \\
			 \hline
			$[5\,000; 6\,000)$ & $24$ \\ 
			\hline 
			 $[6\,000; 7\,000)$ & $29$ \\
			\hline 
			$[7\,000;  8\,000)$ & $18$ \\
			 \hline
			$[8\,000; 9\,000)$ & $9$ \\ 
			\hline 
		\end{tabular}
	\end{center}
	 \choiceTF
		{Cỡ mẫu của mẫu số  liệu bằng $5$}
		{\True Giá  trị đại diện của nhóm thứ $ 4 $ là $7\,500$}
		{\True Số trung bình của mẫu số liệu ghép nhóm bằng $6\,369{,}57$} 
		{\True  Mốt của mẫu số liệu ghép nhóm là $6\,362{,}5$}
	 \loigiai{
		\begin{itemchoice}
			\itemch \textbf{Sai.}\\ 
			Cỡ mẫu của mẫu số  liệu bằng $92$.
			\itemch  \textbf{Đúng. }\\
			Giá trị đại diện của nhóm  thứ $ 4 $ là $7\,500$.
			\itemch \textbf{Đúng. }\\ 
			Số trung bình của mẫu  số liệu ghép nhóm là\\
			$\overline{x}=\dfrac{12\cdot 4\,500+24\cdot  5\,500+29\cdot 6\,500+18\cdot 7\,500+9\cdot 8\,500}{92} \approx 6\,369{,}57$.
			\itemch \textbf{Đúng. } Nhóm chứa mốt  của mẫu số liệu trên là nhóm $[6\,000;7\,000)$.\\
			Do đó $u_m=6\,000$; $n_{m-1}=24$; $n_m=29$; $n_{m+1}=18$;  $u_{m+1}-u_m=7\,000-6\,000=1\,000$.\\
			Mốt của mẫu số liệu ghép  nhóm là
			 $M_o=6\,000+\dfrac{29-24}{(29-24)+(29-18)} \cdot 1\,000=6\,312{,}5$.
		\end{itemchoice}
	 }
\end{ex}

% nguồn: Ok_dungsai_lop11_Lan2
\begin{ex}%[1D5V1-1]%[ma:c34c5428]
	Cho mẫu số liệu ghép nhóm về thống kê chiều cao của $35$ cây bạch đàn trong  rừng, ta có bảng số liệu sau
	\begin{center}
		\begin{tabular}{|c|c|c|c|c|}
			\hline 
			Khoảng chiều cao (m) &  $[6{,}5;7)$ & $[7;7{,}5)$ & $[7{,}5;8)$ & $[8;8{,}5)$ \\
			 \hline
			Số cây & $ 6 $ &  $ 15 $ & $ 11 $ & $ 3 $ \\
			\hline 
		\end{tabular}
	\end{center}
	Xét tính đúng  sai của các mệnh đề sau
	 \choiceTF
		{\True Cỡ mẫu  của mẫu số liệu là $35$}
		{Nhóm chứa trung vị của mẫu số liệu  là $[7{,}5;8{,}0)$}
		{Điểm trung  bình của các học sinh là $7{,}9$ (Kết quả làm tròn đến hàng phần chục)}
		{\True Mốt của mẫu số  liệu là $7{,}35$ (Kết quả làm tròn đến hàng phần trăm)}
	 \loigiai{
		\begin{itemchoice}
			\itemch \textbf{Đúng.  }\\
			Cỡ mẫu là $6+15+11+3=35$. 
			 \itemch \textbf{Sai. }\\
			Nhóm chưa  trunng vị là $[7{,}0;7{,}5)$.
			\itemch \textbf{Sai.  }\\
			Điểm trung bình của các học sinh là\\ $$\overline{x}=\dfrac{6{,}75\cdot  6+7{,}25\cdot 15+7{,}75\cdot 11+8{,}25\cdot 3}{35} \approx 7{,}4.$$
			\itemch \textbf{Đúng.  }\\
			 Nhóm có tần số lớn nhất $[7{,}0; 7{,}5)$.\\
			Ta  có mốt của mẫu số liệu là\\ $$M_o=7{,}0+\dfrac{15-6}{(15-6)+(15-11)}\cdot0{,}5=7{,}35.$$
		\end{itemchoice}
	 }
\end{ex}

% nguồn: Ok_dungsai_lop11_Lan2
\begin{ex}%[1D5V1-1]%[ma:2af16ce2]
	Một nhà nghiên cứu ghi lại thời gian (giờ) sử dụng Facbook của $ 30 $ học sinh trong  $ 02 $ tuần.
	 Kết quả thu được mẫu số liệu như sau:
	\begin{center}
		\begin{tabular}{cccccccccc}
			$21$ & $17$ & $22$ & $18$ & $20$ & $17$ & $15$ & $13$ & $15$ & $20$ \\ $15$ & $12$ & $18$ & $17$ & $25$ & $17$ & $21$ & $15$ & $12$ & $18$ \\ $16$  & $23$ & $14$ & $18$ & $19$ & $13$ & $16$ & $19$ & $18$ & $17$
		\end{tabular}
	\end{center}
	\choiceTF 
		 {Số giờ trung bình của học sinh trong $ 02 $ tuần: $16{,}37$ giờ}
		{\True Tổng hợp kết quả thời gian sử dụng Facbook của  học sinh vào bảng tần số ghép nhóm theo mẫu sau
		\begin{center}
			\begin{tabular}{|c|c|c|c|c|c|}
				\hline 
				Số giờ & $[12;15)$ & $[15;18)$ &  $[18;21)$ & $[21;24)$ & $[24;27)$ \\
				 \hline
				Giá trị đại diện & $13{,} 5$ & $16 {,} 5$ & $18{,}5$ & $21 {,} 5$ &  $24 {,} 5$ \\
				\hline 
				 Số học sinh & $5$ & $12$ & $8$ & $4$ & $1$ \\
				\hline 
		\end{tabular}
	\end{center}}
	{\True  Nhóm chứa mốt của mẫu số liệu ý là nhóm $[15;18)$}
	{\True Mốt  của mẫu số liệu ý bằng $16{,}91$}
	 \loigiai{
		\begin{itemchoice}
			 \itemch \textbf{Sai.} \\
			Tổng số thời gian sử dụng Facbook của 30 học sinh  là: $ 521 $ giờ.\\
			Số giờ trung bình của học sinh trong $  02 $ tuần: $\overline{x}=\dfrac{521}{30}=17{,}37$ giờ.
			\itemch \textbf{Đúng.}  \\
			\begin{center}
				\begin{tabular}{|c|c|c|c|c|c|}
					 \hline
					Số giờ & $[12;15)$ & $[15;18)$  & $[18;21)$ & $[21;24)$ & $[24;27)$ \\
					 \hline
					Giá trị đại diện & $ 13{,}5 $ & $ 16{,}5 $ & $  18{,}5 $ & $ 21{,}5 $ & $ 24{,}5 $ \\
					 \hline
					Số học sinh & $ 5 $ &  $ 12 $ & $ 8 $ & $ 4 $ & $ 1 $ \\
					\hline 
				\end{tabular}
			\end{center}
			\itemch \textbf{Đúng.} \\ 
			{Nhóm chứa mốt của mẫu số liệu trên  là nhóm $[15;18)$}.
			\itemch \textbf{Đúng.} \\ $u_m=15$; $  n_{m-1}=5$; $n_m=12$; $n_{m+1}=8$; $u_{m+1}-u_m=18-15=3{,}0$.
			\[M_{o}=15+\dfrac{12-5}{(12-5)+(12-8)} \cdot  3=16{,}91.
			 \]
		\end{itemchoice}
	 }
\end{ex}

% nguồn: Ok_dungsai_lop11_Lan2
\begin{ex}%[1D5V1-1]%[ma:d64a566c]
	Cho mẫu số liệu ghép nhóm về lương của nhân viên  trong một công ty như sau
	\begin{center}
		\begin{tabular}{|c|c|c|c|c|c|}
			 \hline
			Lương (triệu đồng) & $[9;12)$ & $[12;15)$  & $[15;18)$ & $[18;21)$ & $[21;24)$ \\
			\hline 
			Số nhân viên & $ 6 $  & $ 12 $ & $ 4 $ & $ 2 $& $1 $ \\
			\hline 
		\end{tabular}
	\end{center}
	 \choiceTF
		{\True Giá trị đại diện  của nhóm $[9;12)$ là $10{,}5$}
		{Trung bình  lương các nhân viên là $16{,}5$ triệu đồng}
		{Nhóm chứa trung vị là $[15;18)$} 
		{\True Tứ phân vị thứ ba là $15{,}56$} 
	\loigiai{ 
		\begin{itemchoice}
			 \itemch \textbf{Sai.} \\
			Giá trị đại diện của  nhóm $[9;12)$ là $\dfrac{9+12}{2}=10{,}5$.
			 \itemch \textbf{Sai.} \\
			Trung bình lương các nhân viên là  $\overline{x}=\dfrac{1}{25} (6\cdot 10{,}5+12\cdot 13{,}5+4\cdot 16{,}5+2\cdot 19{,}5+22{,}5)=14{,}1$ triệu đồng.
			 \itemch \textbf{Sai.} \\
			Công ty có $ 25 $ nhân sự. Vì $x_{13} \in [12;15)$ nên nhóm này chứa trung vị. 
			\itemch \textbf{Sai.} \\ 
			Vì $x_{19} \in [15;18)$ nên ta có $Q_3=15+\dfrac{\dfrac{3\cdot 25}{4}-18}{4} \cdot  3\approx 15{,}56$.
		\end{itemchoice}
	} 
\end{ex}

% nguồn: Ok_dungsai_lop11_Lan2
\begin{ex}%[1D5V1-1]%[ma:00710c30]
	Chiều cao  của $ 35 $ cây bạch đàn sinh trưởng của $ 12 $ tháng được ghi lại ở bảng sau (đơn vị: mét)
	\begin{center}
		\begin{tabular}{|c|c|c|c|c|c|c|c|c|c|c|c|}
			 \hline
			$6{,}6$ & $7{,}5$ & $8{,}2$ & $8{,}2$ & $7{,}8$ & $7{,}9$  & $9{,}0$ & $8{,}9$ & $8{,}2$ \\
			\hline 
			$7{,}2$ & $7{,}5$ & $8{,}3$ & $7{,}4$ & $8{,}7$ & $7{,}7$ & $7{,}0$ & $9{,}4$ & $8{,}7$  \\
			\hline 
			$8{,}0$ & $7{,}7$ & $7{,}8$ & $8{,}3$  & $8{,}6$ & $8{,}1$ & $8{,}1$ & $9{,}5$ & $6{,}9$ \\
			 \hline
			$8{,}0$ & $7{,}6$ & $7{,}9$ & $7{,}3$ & $8{,}5$  & $8{,}4$ & $8{,}0$ & $8{,}8$ & \\
			 \hline
		\end{tabular}
	\end{center}
	\choiceTF 
		{\True Chiều cao sinh trưởng trung  bình của mỗi tháng của cây bạch đàn là: $8{,}05\; $m}
		{\True  Tổng hợp lại kết quả chiều cao của cây bạch đàn vào bảng tần số ghép nhóm theo mẫu sau
		\begin{tabular}{|c|c|c|c|c|c|c|}
			 \hline
			Chiều cao (m) & $[6{,}5;7{,}0)$ & $[7{,}0;7{,}5)$ & $[7{,}5;8)$  & $[8;8{,}5)$ & $[8{,}5;9{,}0)$ & $[9{,}0;9{,}5)$ \\
			\hline 
			Số cây & $2$ & $4$ & $9$ & $11$ & $6$ & $3$  \\
			\hline 
	\end{tabular}}
	{\True Ước  lượng chiều cao trung bình mỗi tháng của cây bạch đàn từ bảng tần số ghép nhóm ở ý là $ 8{,}09$ ( m)}
	{Chiều cao tăng trưởng của cây bạch đàn được $ 8{,}18 $  m là cao nhất}
	 \loigiai{
		\begin{itemchoice}
			\itemch \textbf{Đúng.}  \\
			{Chiều cao sinh  trưởng trung bình của mỗi tháng của cây bạch đàn là $8{,}05\; $ m}
			\itemch \textbf{Đúng.}\\ \begin{tabular}{|c|c|c|c|c|c|c|}
			\hline 
			Chiều cao $(m)$ & $[6{,}5;7{,}0)$ & $[7{,}0;7{,}5)$ & $[7{,}5;8)$ & $[8;8{,}5)$ &  $[8{,}5;9{,}0)$ & $[9{,}0;9{,}5)$ \\
			\hline 
			Số cây & $ 2  $ & $ 4 $ & $ 9 $ & $ 11 $ & $ 6 $ & $ 3 $ \\
			\hline 
		\end{tabular}
		\itemch \textbf{Đúng.}\\ \begin{tabular}{|c|c|c|c|c|c|c|}
		 \hline
		Chiều cao (m) &  $[6{,}5;7{,}0)$ & $[7{,}0;7{,}5)$ & $[7{,}5;8)$ & $[8;8{,}5)$ & $[8{,}5;9{,}0)$ & $[9{,}0;9{,}5)$ \\
		\hline 
		Giá trị đại diện & $6 {,} 75$  & $7{,}25$ & $7{,}75$ & $8{,}25$ & $8{,}75$ & $9 {,} 25$ \\
		 \hline
		Số lần & $2$  & $4$ & $9$ & $11$ & $6$ & $3$ \\
		\hline 
	\end{tabular}
	Chiều cao trung  bình mỗi tháng sinh trưởng xấp xỉ bằng
	\[\dfrac{6{,}75\cdot  2+7{,}25\cdot4+7{,}75\cdot9+8{,}25\cdot11+8{,}75\cdot6+9{,}25\cdot3}{35} \approx 8{,}09 \text{(m).}
	\] 
	\itemch \textbf{Sai.}  \\
	Nhóm chứa mốt của mẫu  số liệu trên là $[8;8{,}5)$.
	Do đó\\ $u_m=8$; $n_{m-1}=9$; $n_{m+1}=6$; $u_{m+1}-u_m=8{,}5-8=0{,}5$.\\ 
	Mốt của mẫu số liệu ghép nhóm là $M_o=8+\dfrac{11-9}{(11-9)+(11-6)}  \cdot 0{,}5\approx 8{,}14$.\\
	Vậy chiều cao tăng trưởng của cây bạch đàn được $ 8{,}14 $ m là cao nhất. 
	\end{itemchoice}
 }
\end{ex}

% nguồn: Ok_dungsai_lop11_Lan2
\begin{ex}%[1D5V1-1]%[ma:9c1a9ee3]
	Cho mẫu số  liệu điểm môn Toán của một nhóm học sinh như sau
	\begin{center}
		\begin{tabular}{|c|c|c|c|c|}
			\hline 
			 Điểm & $[6;7)$ & $[7;8)$ & $[8;9)$ & $[9;10)$ \\
			\hline 
			Số học sinh & $8$ & $7$ & $10$ &  $5$ \\
			\hline 
		\end{tabular}
	\end{center}
	 \choiceTF
		{\True Mẫu số liệu đã cho  là mẫu số liệu ghép nhóm}
		{\True Cỡ mẫu của mẫu số liệu là $30$} 
		 {\True Điểm trung bình của các học sinh là $7{,}9$}
		 {Mốt của mẫu số liệu là $10$}
	\loigiai{ 
		Ta có bảng  sau
		\begin{center}
			\begin{tabular}{|c|c|c|c|c|}
				\hline 
				 Điểm & $[6;7)$ & $[7;8)$ & $[8;9)$ & $[9;10)$ \\
				 \hline
				Giá trị đại diện & $ 6{,}5 $ & $ 7{,}5 $ & $ 8{,}5 $ & $ 9{,}5 $  \\
				\hline 
				Số học sinh & $ 8  $ & $ 7 $ & $ 10 $ &$ 5 $ \\
				\hline 
			\end{tabular}
		\end{center}
		\begin{itemchoice}
			\itemch \textbf{Đúng.}  \\
			{Theo định nghĩa mẫu số liệu ghép nhóm}. 
			\itemch \textbf{Đúng.}\\ 
			{Cỡ mẫu  là $8+7+10+5=30$}.
			\itemch \textbf{Đúng.} \\ 
			{Điểm trung bình của các học sinh là $\overline{x}=\dfrac{6{,}58+7{,}5\cdot 7+8{,}5\cdot  10+9{,}5\cdot 5}{30}=7{,}9$}.
			\itemch  \textbf{Sai.} \\
			Nhóm  chứa Mốt là $[8;9)$.\\
			Mốt của mẫu số liệu  là $M_o=8+\dfrac{10-7}{10-7+10-5} (9-8)=8{,}375$.
		\end{itemchoice}
	 }
\end{ex}

% nguồn: Ok_dungsai_lop11_Lan2
\begin{ex}%[1D5V1-1]%[ma:3bd0c8a6]
	Người ta tiến hành phỏng vấn $ 30 $ người về một bộ phim mới chiếu trên truyền hình. Người điều tra yêu cầu cho điểm bộ phim (thang điểm là $ 100 $). Kết quả được trình bày  trong bảng phân bố tần số ghép lớp sau đây
	\begin{center}
		\begin{tabular}{|c|c|c|c|c|c|}
			\hline 
			 Số điểm & $[50;60)$ & $[60;70)$ & $[70;80)$ & $[80;90)$ & $[90;100)$ \\
			 \hline
			Số người & $2$ &  $6$ & $10$ & $8$ & $4$ \\
			 \hline
		\end{tabular}
	\end{center}
	\choiceTF 
		{\True Ước  lượng số trung bình của mẫu ghép là $77$}
		{\True Giá trị đại diện  của nhóm $[90;100)$ là $ 95 $}
		{Nhóm chứa mốt của mẫu số liệu  trên là nhóm $[80;90)$}
		{Mốt  của mẫu số liệu là $74{,}67$}
	 \loigiai{
		\begin{center}
			\begin{tabular}{|c|c|c|c|c|c|}
				\hline 
				Số điểm & {$[50;60)$} & $[60;70)$ & {$[70;80)$} & {$[80;90)$} & $[90;100)$ \\ 
				 \hline
				Giá trị đại diện & $ {55} $ & $ 65 $ & $ {75} $ & $ {85} $ & $ 95 $  \\
				 \hline
				Số người & $ {2} $ & $ 6 $ & $ {10} $ & $  {8} $ & $ 4 $ \\
				\hline 
			\end{tabular}
		\end{center}
		\begin{itemchoice}
			\itemch  \textbf{Đúng.}\\
			Ước lượng số trung bình của mẫu  ghép là
			\[\overline{x}=\dfrac{55\cdot 2+65\cdot 6+75\cdot 10+85\cdot 8+95 \cdot4}{30}=77. 
			 \]
			 \itemch \textbf{Đúng.}\\
			 Giá trị đại diện của nhóm $[90;100)$ là $ 95 $.
			 \itemch \textbf{Sai.} \\
			Nhóm chứa mốt của mẫu số liệu  trên là nhóm $[70;80)$.
			\itemch \textbf{Sai.}\\ $u_m=70$; $n_{m-1}=6$;  $n_m=10$; $n_{m+1}=8$; $u_{m+1}-u_m=80-70=10$.\\
			\[M_{o}=70+\dfrac{10-6}{(10-6)+(10-8)} \cdot 10=76{,}67.\; 
			 \]
		\end{itemchoice}
	 }
\end{ex}

% nguồn: Ok_dungsai_lop11_Lan2
\begin{ex}%[1D5V1-1]%[ma:6babb215]
	Kết quả khảo sát cân nặng  của $ 25 $ quả cam ở mỗi lô hàng $ A $, $ B $ được cho ở bảng sau
	\begin{center}
		\begin{tabular}{|c|c|c|c|c|c|}
			\hline 
			Cân nặng (gam) & $[150;155)$ & $[155;160)$ & $[160;165)$ & $[165;170)$ & $[170;175)$  \\
			\hline 
			ố quả cam ở lô hàng $A$ & $2$ & $6$ & $12$ & $4$  & $1$ \\
			\hline 
			Số  quả cam ở lô hàng $B$ & $1$ & $3$ & $7$ & $10$ & $4$ \\
			\hline 
		\end{tabular}
	\end{center}
	\choiceTF 
		{\True Giá trị đại diện nhóm $[150;155)$ bằng $ 152{,}5 $} 
		{Cân nặng  trung bình của mỗi quả cam ở lô $A$ là $163{,}7\; $(gam)}
		{Cân nặng trung bình của mỗi quả cam  ở lô $B$ là $162{,}1\; $(gam)}
		{\True Theo số trung bình  thì cam ở lô hàng $B$ nặng hơn cam ở lô hàng $A$}
	 \loigiai{
		Bảng thống kê  số lượng cam theo giá trị đại diện
		\begin{center}
			\begin{tabular}{|c|c|c|c|c|c|}
				\hline 
				Cân nặng đại diện (gam) & $ 152{,}5 $ & $ 157{,}5  $ & $ 162{,}5 $ & $ 167{,}5 $ & $ 172{,}5 $ \\
				 \hline
				Số quả cam ở lô hàng $A$ & $2$ & $6$ & $12$ & $4$ & $1$ \\ 
				 \hline
				Số  quả cam ở lô hàng $B$ & $1$ & $3$ & $7$ & $10$ & $4$ \\
				 \hline
			\end{tabular}
		\end{center}
		\begin{itemchoice}
			\itemch \textbf{Đúng. }Giá trị đại diện nhóm $[150;155)$ bằng $ 152{,}5  $\\
			\itemch \textbf{Sai.} Cân nặng trung bình  của mỗi quả cam ở lô $A$ là\\
			\[\overline{x}_A=\dfrac{152{,}5\cdot 2+157{,}5\cdot 6+162{,}5\cdot 12+167{,}5\cdot 4+172{,}5\cdot 1}{25}=161{,}7\; \text{(gam).} 
			\] 
			\itemch \textbf{Sai.} Cân nặng trung bình của  mỗi quả cam ở lô $B$ là\\
			\[\overline{x}_B=\dfrac{152{,}5\cdot 1+157{,}5\cdot  3+162{,}5\cdot 7+167{,}5\cdot 10+172{,}5\cdot 4}{25}=165{,}1\; \text{(gam).}
			 \]
			\itemch  \textbf{Đúng. } Ta thấy $\overline{x}_A < \overline{x}_B$. Vậy nếu so sánh theo số trung bình thì cam ở lô hàng $B$ nặng hơn cam ở lô hàng $A$.
		\end{itemchoice}
	} 
\end{ex}

% nguồn: Ok_dungsai_lop11_Lan2
\begin{ex}%[1D5V1-1]%[ma:affff56d]
	Một thư viện thống kê số lượng sách  được mượn mỗi ngày trong một tháng ở bảng sau
	\begin{center}
		\begin{tabular}{|c|c|c|c|c|c|}
			\hline 
			Số sách & $ [16;20) $ &$ [21;25) $  & $ [26;30) $ & $ [31;35) $ & $ [36;40) $ \\
			 \hline
			 Số ngày & $ 4 $ & $ 5 $ & $ 11 $ & $ 7 $ & $ 3 $ \\
			 \hline
		\end{tabular}
	\end{center}
	\choiceTF 
		{\True Giá trị đại diện của nhóm $[16;20]$ là $ 18 $} 
		{Nhóm chứa mốt của mẫu  số liệu là $[36;40]$}
		{\True Số trung bình của mẫu số  liệu ghép nhóm trên bằng $28$}
		{\True Mốt của mẫu số liệu ghép nhóm  trên bằng $28{,}4$}
	\loigiai{ 
		\begin{itemchoice}
			 \itemch \textbf{Đúng.}\\
			{Giá trị đại diện của nhóm  $[16;20]$ là $\dfrac{1}{2} (16+20)=18$.}
			\itemch  \textbf{Sai.}\\
			{Nhóm chứa mốt của mẫu số liệu  là $[26;30]$.}
			\itemch \textbf{Đúng.}\\ 
			 {Số trung bình bằng: \[(18\cdot 4+23\cdot 5+28\cdot 11+33\cdot 7+38\cdot 3):30=28.\] }
			\itemch \textbf{Đúng.}\\ Nhóm chứa  mốt là $[26;30]$.\\
			Ta có $u_m=26,n_{m-1}=5,n_m=11,n_{m+1}=7,u_{m+1}-u_m=30-26=4$.\\ 
			Mốt của mẫu  số liệu là $M_o=26+\dfrac{11-5}{(11-5)+(11-7)}\cdot 4=28{,}4$.
		\end{itemchoice}
	 }
\end{ex}

% nguồn: Ok_dungsai_lop11_Lan2
\begin{ex}%[1D5V1-1]%[ma:544ab7c2]
	Số cuộc  điện thoại một người thực hiện mỗi ngày trong $ 30 $ ngày được lựa chọn ngẫu nhiên được thống kê trong bảng sau
	\begin{center}
		\begin{tabular}{|c|c|c|c|c|c|}
			\hline 
			Số cuộc gọi & $[2{,}5;5{,}5)$ & $[5{,}5;8{,}5)$ & $[8{,}5;11{,}5)$ & $[11{,}5;14{,}5)$ &  $[14{,}5;17{,}5)$ \\
			 \hline
			Số ngày & 5 &  13 & 7 & 3 & 2 \\
			 \hline
		\end{tabular}
	\end{center}
	\choiceTF 
		 {\True Số cuộc gọi trung bình mỗi ngày là: $8{,}1$}
		{Nhóm chứa  mốt là $[5.5;8{,}5)$}
		{\True Mốt  của mẫu số liệu ghép nhóm là $7{,}21$}
		{Người đó thực hiện tối đa khoảng $  8 $ cuộc gọi mỗi ngày}
	\loigiai{ 
		 Ta viết lại bảng tần số ghép nhóm theo giá trị đại diện là
		\begin{center}
			\begin{tabular}{|c|c|c|c|c|c|}
				\hline 
				Số cuộc gọi  & $[2{,}5;5{,}5)$ & $[5{,}5;8{,}5)$ & $[8{,}5;11{,}5)$ & $[11{,}5;14{,}5)$ & $[14{,}5;17{,}5)$ \\
				\hline 
				Giá trị đại diện & $ 4 $ & $ 7 $ & $ 10 $ & $ 13  $ & $ 16 $ \\
				\hline 
				Số ngày &$ 5 $ & $ 13 $ &$ 7 $  & $ 3 $ & $ 2 $ \\
				\hline 
			\end{tabular}
		\end{center}
		\begin{itemchoice}
			\itemch  \textbf{Sai.} \\Số cuộc gọi trung bình mỗi ngày là
			\[\overline{x}=\dfrac{4\cdot5+7\cdot 13+10 \cdot 7+13\cdot  3+16\cdot 2}{30}=8{,}4.
			\] 
			\itemch \textbf{Đúng.}  \\
			 Nhóm chứa mốt là $[5{,}5;8{,}5)$.
			\itemch \textbf{Đúng.}\\ $u_m=5{,}5$; $u_{m+1}=8{,}5\Rightarrow u_{m+1}-u_m=3$;  $n_m=13$; $n_{m-1}=5$; $n_{m+1}=7$.
			\itemch \textbf{Sai.} \\Mốt của mẫu  số liệu ghép nhóm là\\
			\[M_o=5{,}5+\dfrac{13-5}{(13-5)+(13-7)} \cdot 3=\dfrac{101}{14} \approx  7{,}21.
			\] 
		\end{itemchoice}
	} 
\end{ex}

% nguồn: Ok_dungsai_lop11_Lan2
\begin{ex}%[1D5V1-1]%[ma:31c96322]
	Khảo sát chiều cao của học  sinh lớp 11C1, thu được bảng số liệu ghép nhóm như sau
	\begin{center}
		\begin{tabular}{|c|c|c|c|c|c|}
			 \hline
			Chiều cao (cm) & $[150;155)$ & $[155;160)$ & $[160;165)$ & $[165;170)$ & $[170;175)$ \\ 
			\hline 
			Số học sinh & $ 7 $ & $ 15 $ & $ 12 $  & $ 8 $& $3 $ \\
			\hline 
		\end{tabular}
	\end{center}
	\choiceTF 
		{Lớp 11C1  có $ 43 $ học sinh}
		{\True Chiều cao trung bình của học sinh 11C1  thuộc khoảng $(161; 162)$}
		{Số học sinh có chiều cao  không dưới $ 160 $ cm là $ 22 $ học sinh}
		{\True Trung vị của mẫu  số liệu ghép nhóm trên là $\dfrac{961}{6}$ cm}
	\loigiai{ 
		\begin{itemchoice}
			\itemch \textbf{Sai.} \\ 
			Số học sinh lớp 11C1 là $7+15+12+8+3=45$  học sinh.
			\itemch \textbf{Đúng.}  \\
			Chiều cao trung bình của học sinh lớp 11C1 là\\ $\overline{x}=\dfrac{1}{45}  (7\cdot 152{,}5+15\cdot 157{,}5+12\cdot 162{,}5+4\cdot 167{,}5+3\cdot 172{,}5)=160{,}8$ cm.
			 \itemch \textbf{Sai.}\\
			Số học sinh có  chiều cao không dưới $ 160 $ cm là $12+8+3=23$ học sinh.
			\itemch  \textbf{Đúng.}\\
			Vì $x_{13} \in [155;160)$ nên ta  có $M_e=155+\dfrac{\dfrac{45}{2}-7}{15} \cdot (160-155)=\dfrac{961}{6}$ cm.
		\end{itemchoice}
	} 
\end{ex}

% nguồn: Ok_dungsai_lop11_Lan2
\begin{ex}%[1D5V1-1]%[ma:a357d2b4]
	Thâm niên giảng dạy của một số  giáo viên trường THPT được ghi lại ở bảng sau
	\begin{center}
		\begin{tabular}{|c|c|c|c|c|c|}
			 \hline
			 Thâm niên (Số năm) & $[1;5)$ & $[5;10)$ & $[10;15)$ & $[15;20)$ & $[20;25)$ \\
			 \hline
			Số giáo viên & $ 4 $ & $ 12 $ & $ 16 $  & $ 8 $ & $ 3 $ \\
			\hline 
		\end{tabular}
	\end{center}
	\choiceTF 
		{\True Cỡ mẫu của mẫu  số liệu bằng $43$}
		{\True  Số trung bình của mẫu ghép nhóm là $11{,}84$}
		{\True Nhóm chứa mốt của mẫu số  liệu trên là nhóm $[10;15)$}
		 {Mốt của mẫu số liệu ghép nhóm bằng $11{,}74$}
	\loigiai{ 
		\begin{center}
			\begin{tabular}{|c|c|c|c|c|c|}
				 \hline
				Thâm niên (Số năm) & {$[1;5)$}  & $[5;10)$ & $[10;15)$ & {$[15;20)$} & $[20;25)$ \\
				 \hline
				Giá trị đại diện & $ 3 $ & $ 7{,}5 $ & $ 12{,}5 $ &  $ 17{,}5 $ & $ 22{,}5 $ \\
				\hline 
				Số giáo viên  & {$ 4 $} & $ 12 $ & $ 16 $ & {$ 8 $} & $ 3 $ \\
				 \hline
			\end{tabular}
		\end{center}
		\begin{itemchoice}
			\itemch  \textbf{Đúng.} \\
			Cỡ mẫu của mẫu số liệu bằng $  4+12+16+8+3=43 $.
			 \itemch \textbf{Đúng.} \\
			Ước  lượng số trung bình của mẫu ghép là
			\[\overline{x}=\dfrac{3\cdot4+7{,}5\cdot12+12{,}5\cdot16+17{,}5\cdot8+22{,}5\cdot3}{43} \approx 11{,}84. 
			\] 
			\itemch  \textbf{Đúng.} \\
			Nhóm chứa mốt của mẫu số  liệu trên là nhóm $[10;15)$.
			\itemch \textbf{Sai.}\\ $u_m=10$; $ n_{m-1}=12$;  $ n_m=16$; $ n_{m+1}=8$; $n_{m+1}-u_m=15-10=5$.
			\[M_{^\circ}=10+\dfrac{16-12}{(16-12)+(16-8)}  \cdot 5=11{,}66.
			 \]
		\end{itemchoice}
	 }
\end{ex}

% nguồn: Ok_dungsai_lop11_Lan2
\begin{ex}%[1D5V1-1]%[ma:5956f749]
	Anh An ghi lại cự li $ 30 $ lần ném lao của mình  ở bảng sau (đơn vị: mét)
	\begin{center}
		\begin{tabular}{|c|c|c|c|c|c|c|c|c|c|}
			 \hline
			$ 72{,}1 $ & $ 72{,}9 $ & $ 70{,}2 $ & $ 70{,}9 $ & $ 72{,}2 $ & $ 71{,}5 $ & $ 72{,}5 $ & $ 69{,}3 $ &  $ 72{,}3 $ & $ 69{,}7 $ \\ \hline
			$ 72{,}3 $ & $ 71{,}5 $ & $ 71{,}2 $ & $ 69{,}8 $ & $ 72{,}3 $ & $ 71{,}1  $ & $ 69{,}5 $ & $ 72{,}2 $ & $ 71{,}9 $ & $ 73{,}1 $ \\ \hline
			$ 71{,}6 $ & $ 71{,}3 $ & $ 72{,}2 $ & $ 71{,}8 $ & $ 70{,}8 $ & $ 72{,}2 $  & $ 72{,}2 $ & $ 72{,}9 $ & $ 72{,}7 $ & $ 70{,}7 $ \\ \hline
		\end{tabular}
	\end{center}
	\choiceTF 
		{Cự li trung bình của mỗi  lần ném là $ 72{,}3 $ (m)}
		{\True Số lần ném của anh An cự li $ [70{,}8;71{,}8] $ là $ 10  $}
		{\True Tổng hợp lại kết quả ném của anh An vào bảng tần số ghép nhóm ta được  bảng sau
		\begin{tabular}{|c|c|c|c|c|c|}
			 \hline
			Cự li (m) & $[69{,}2;70)$  & $[70;70{,}8)$ & $[70{,}8;71{,}6)$ & $[71{,}6;72{,}4)$& $[72{,}4;73{,}2)$ \\
			 \hline
			Số lần & $ 4 $ & $ 2  $ & $ 9 $ & $ 10 $ & $ 5 $ \\
			 \hline
	\end{tabular}\\
	Khi đó cự li trung bình mỗi lần ném từ bảng tần số ghép nhóm trên là  $ 71{,}5 $ m}
	{Khả năng anh An ném được cao nhất là  $ 72 $ m}
	 \loigiai{
		\begin{itemchoice}
			\itemch  \textbf{Sai.} \\
			Cự li trung bình của mỗi lần ném  là $ 71{,}6 $ (m).
			 \itemch \textbf{Đúng. } \\
			Số lần ném của anh An cự li $ [70{,}8;71{,}8] $ là $  10 $.
			\itemch \textbf{Đúng.  } \\
			Cự li trung  bình mỗi lần ném xấp xỉ bằng\\
			$\dfrac{69{,}6\cdot 4+70{,}4\cdot  2+71{,}2\cdot 9+72\cdot 10+72{,}8\cdot 5}{30}=71{,}5(\; m)$.
			\itemch \textbf{Sai.}\\ $u_m=71{,}6$; $n_{m-1}=9$; $n_{m+1}=5$; $u_{m+1}-u_m=72{,}4-71{,}6=0{,}8$.\\ 
			Mốt của  mẫu số liệu ghép nhóm là
			\[M_o=71{,}6+\dfrac{10-9}{(10-9)+(10-5)} \cdot  0{,}8=71{,}7\text{(m).}
			 \]
			Vậy khả năng anh  An ném được $ 71{,}7 $ m là cao nhất.
		\end{itemchoice}
	 }
\end{ex}

% nguồn: Ok_dungsai_lop11_Lan2
\begin{ex}%[1D5V1-1]%[ma:2978d11f]
	Thống kê điểm trung bình môn Toán của một số học sinh  được cho ở bảng sau
	\begin{center}
		\begin{tabular}{|c|c|c|c|c|c|c|c|}
			\hline 
			Doanh thu & $ [6{,}5;7) $ &$ [7;7{,}5) $ & $ [7{,}5;8) $ & $ [8;8{,}5) $ & $ [8{,}5;9) $ & $ [9;9{,}5) $ &  $ [9{,}5;10) $ \\
			\hline 
			Số ngày & $ 8  $ & $ 10 $ & $ 15 $ & $ 25 $ & $ 13 $ & $ 7 $ & $ 3 $ \\
			\hline 
		\end{tabular}
	\end{center}
	 \choiceTF
		 {\True Cỡ mẫu $n=81$}
		{\True Số trung bình của mẫu số liệu trên thuộc khoảng $[8;8{,}5)$} 
		{Số  trung vị của mẫu số liệu trên thuộc khoảng $[7{,}5;8)$}
		{Mốt của mẫu  số liệu trên thuộc khoảng $[7{,}5;8)$}
	 \loigiai{
		\begin{itemchoice}
			 \itemch \textbf{Đúng.}\\
			Cỡ mẫu  $n=8+10+15+25+13+7+3=81$.
			\itemch \textbf{Đúng.} 
			\begin{center}
				\begin{tabular}{|c|c|c|c|c|c|c|c|}
					\hline 
					Doanh thu  & $ [6{,}5;7) $ &$ [7;7{,}5) $ & $ [7{,}5;8) $ & $ [8;8{,}5) $ & $ [8{,}5;9) $ & $ [9;9{,}5) $ & $ [9{,}5;10) $ \\
					 \hline
					Điểm đại diện & $ 6{,}75 $ & $ 7{,}25 $ & $ 7{,}75 $ & $ 8{,}25 $ & $ 8{,}75 $ & $ 9{,}25  $ & $ 9{,}75 $ \\
					 \hline
					Số ngày & $ 8 $ & $ 10 $ & $ 15 $ & $ 25 $ & $ 13  $ & $ 7 $ & $ 3 $ \\
					 \hline
				\end{tabular}
			\end{center}
			Số trung bình bằng $$(6{,}75\cdot 8+7{,}25\cdot 10+7{,}75\cdot 15+8{,}25 \cdot 25+8{,}75  \cdot 13+9{,}25 \cdot7+9{,}75 \cdot 3):81\approx 8{,}1.$$
			\itemch  \textbf{Sai.}\\
			Gọi $x_1;x_2;\cdots;x_{81}$ là điểm  trung bình của $ 81 $ học sinh theo thứ tự không giảm.\\
			Trung vị của mẫu số liệu là $x_{41}$.\\ 
			Ta  có $n=81$; $n_m=25$; $C=33$; $u_m=8$; $u_{m+1}=8{,}5$.\\
			Trung vị của mẫu số liệu là  $M_e=8+\dfrac{\dfrac{81}{2}-33}{25}\cdot(8{,}5-8)=8,15$.
			\itemch \textbf{Sai.}\\ 
			Nhóm chứa  mốt là $[8;8{,}5)$.\\
			 Ta có $u_m=8$; $n_{m-1}=15$; $n_m=25$; $n_{m+1}=13$; $u_{m+1}-u_m=0{,}5$.\\
			Mốt của mẫu số liệu ghép nhóm là $M_o=8+\dfrac{25-15}{(25-15)+(25-13)}\cdot 0{,}5=\dfrac{181}{22} \approx  8{,}22$.
		\end{itemchoice}
	} 
\end{ex}

% nguồn: Ok_dungsai_lop11_Lan2
\begin{ex}%[1D5V1-1]%[ma:578d9532]
	Một công ty xây dựng  khảo sát khách hàng xem họ có nhu cầu mua nhà ở mức giá nào. Kết quả khỏa sát được lưu lại ở bảng sau
	\begin{center}
		\begin{tabular}{|c|c|c|c|c|c|}
			 \hline
			Mức giá (triệu đồng/m$^2 $) & $[10;14)$ & $[14;18)$ & $[18;22)$ & {$[22;26)$}& $[26;30)$  \\
			\hline 
			Số khách hàng & $ 54 $ & $ 78 $ & $ 120 $ &  $ 45 $ & $ 12 $ \\
			 \hline
		\end{tabular}
	\end{center}
	\choiceTF 
		{\True Độ dài của nhóm  $[10;14)$ bằng $ 4 $}
		{\True Giá trị đại diện  của nhóm $ [10{,}14) $ là $ 12 $}
		{Số trung bình của mẫu số liệu ghép  nhóm bằng $19{,}52$}
		{\True Mốt của  mẫu số liệu ghép nhóm là $19{,}44$}
	\loigiai{ 
		\begin{itemchoice}
			\itemch  \textbf{Đúng. }\\
			Độ dài  của nhóm $[10;14)$ bằng $ 4 $.
			\itemch \textbf{Đúng. } Giá  trị đại diện của nhóm $ [10{,}14) $ là $ 12 $.
			\itemch \textbf{Sai. } \\Số  trung bình của mẫu số liệu ghép nhóm là\\
			$\overline{x}=\dfrac{12\cdot 54+16\cdot 78+20\cdot  120+24\cdot 45+28\cdot 12}{54+78+120+45+12}=18{,}49$.
			\itemch \textbf{Đúng.  } \\Nhóm chứa mốt của mẫu số liệu trên là nhóm $[18;22)$.\\
			Mốt của mẫu số liệu ghép nhóm  là
			\[M_o=18+\dfrac{120-78}{(120-78)+(120-45)} \cdot  (22-18)=19{.}44.\]
		\end{itemchoice}
	} 
\end{ex}

% nguồn: Ok_dungsai_lop11_Lan2
\begin{ex}%[1D5V1-1]%[ma:dee6a717]
	Số người đến thư viện đọc sách trong $30$ ngày của tháng $ 4 $  ở một thư viện được thống kê như bảng dưới
	\begin{center}
		\begin{tabular}{|c|c|c|c|c|c|}
			 \hline
			Lượng người & $[24;30)$ & $[30;36)$ & $[36;42)$ & {$[42;48)$}& $[48;54)$  \\
			 \hline
			Tần số & $ 3 $ & $ 7 $ & $  10 $ & $ 6 $ & $ 4 $ \\
			\hline 
		\end{tabular}
	\end{center}
	 \choiceTF
		{Độ dài của  mỗi nhóm là $5$}
		{\True Giá trị đại diện của nhóm $[30;36)$ là  $33$}
		{Số người đến đọc trung bình mỗi ngày  của thư viện (làm tròn đến hàng đơn vị) là $40$}
		{\True Mốt của mẫu số liệu (làm tròn đến hàng đơn vị) là  $36$}
	 \loigiai{
		\begin{itemchoice}
			\itemch \textbf{Sai.}  \\
			Độ dài của mỗi nhóm là  $30-24=6$.
			 \itemch \textbf{Đúng. }\\
			Giá trị đại diện của nhóm  $[30;36)$ là $\dfrac{30+36}{2}=33$.
			 \itemch \textbf{Sai. }\\
			 Số người đến đọc trung bình mỗi ngày của thư viện là
			\[\dfrac{3\cdot27+7\cdot33+10\cdot39+6\cdot45+4\cdot51}{3+7+10+8+2} \approx  39.
			 \]
			\itemch \textbf{Đúng.  } \\
			Nhóm có tần số lớn nhất là nhóm $[36;42)$. Vậy mốt của mẫu số liệu là 
			\[M_o=36+\dfrac{10-7}{2\cdot10-7-6}\cdot6\approx 36. 
			\] 
		\end{itemchoice}
	} 
\end{ex}

% nguồn: Ok_dungsai_lop11_Lan2
\begin{ex}%[1D5V1-1]%[ma:540a8a0a]
	Doanh thu bán hàng trong $ 20 $ ngày được lựa chọn ngẫu  nhiên của một cửa hàng được ghi lại ở bảng sau (đơn vị: triệu đồng).
	\begin{center}
		\begin{tabular}{|c|c|c|c|c|c|}
			\hline 
			Doanh thu & $ [5;7) $ &$ [7;9) $  & $ [9;11) $ & $ [11;13) $ & $ [13;15) $ \\
			\hline 
			Số ngày & $ 2 $ & $  5 $ & $ 8 $ & $ 4 $ & $ 1 $ \\
			 \hline
		\end{tabular}
	\end{center}
	\choiceTF 
		{Nhóm  chứa mốt của mẫu số liệu trên là $[13;15)$}
		{\True Số trung bình của mẫu số liệu trên  thuộc khoảng $[9;11)$}
		{\True Mốt của mẫu số  liệu trên thuộc khoảng $[9;11)$}
		{Số trung vị của mẫu số liệu trên thuộc khoảng $[11;13)$} 
	 \loigiai{
		\begin{itemchoice}
			\itemch \textbf{Sai.}\\ 
			Nhóm chứa mốt  của mẫu số liệu trên là $[9;11)$.
			 \itemch \textbf{Đúng.}\\
			\begin{center}
				\begin{tabular}{|c|c|c|c|c|c|}
					 \hline
					Doanh thu & $ [5;7) $ &$ [7;9) $ & $ [9;11) $ & $ [11;13) $ & $ [13;15) $  \\
					 \hline
					Giá trị đại diện & $ 6 $ & $ 8 $ & $ 10 $  & $ 12 $ & $ 14 $ \\
					\hline 
					Số ngày & $ 2 $ &  $ 5 $ & $ 8 $ & $ 4 $ & $ 1 $ \\
					 \hline
				\end{tabular}
			\end{center}
			Số trung bình  của mẫu số liệu trên là $\dfrac{6\cdot2+8\cdot5+10\cdot8+12\cdot4+14\cdot1}{20}=9{,}7\in [9;11)$.
			\itemch \textbf{Đúng.}\\ 
			Nhóm chứa mốt  là $[9;11)$.
			Ta có  $u_m=9;n_{m-1}=5$; $n_m=8$; $n_{m+1}=4$; $u_{m+1}-u_m=2$.\\
			Mốt của mẫu số liệu ghép nhóm là $M_o=9+\dfrac{8-5}{(8-5)+(8-4)}\cdot 2=\dfrac{69}{7}  \approx 9{,}85$.
			\itemch \textbf{Sai.} Gọi $x_1;x_2;\ldots;x_{20}$ là doanh thu bán hàng trong 20 ngày theo thứ  tự không giảm.\\
			Trung vị của mẫu số liệu là $\dfrac{1}{2} \bigl(x_{10}+x_{11} \bigr)\in  [9;11)$.\\
			Ta có $n=20$; $n_m=8$; $C=7$; $u_m=9$;  $u_{m+1}=11$.\\
			Trung vị của mẫu số liệu là  $M_e=9+\dfrac{\dfrac{20}{2}-7}{8}\cdot(11-9)=9{,}75\in [9;11)$.
		\end{itemchoice}
	} 
\end{ex}

% nguồn: Ok_dungsai_lop11_Lan2
\begin{ex}%[1D5V1-1]%[ma:c34bed43]
	Khảo sát thời gian tập thể dục trong ngày của một số  học sinh khối $ 11 $ thu được mẫu số liệu ghép nhóm sau
	\begin{center}
		\begin{tabular}{|c|c|c|c|c|c|}
			\hline 
			Thời gian (phút) & $[0;20)$ & $[20;40)$ & $[40;60)$ &  {$[60;80)$}& $[80;100)$ \\
			 \hline
			Số học sinh & $ 5 $ & $ 9 $ & $ 12 $ & $ 10 $ & $ 6 $ \\ 
			 \hline
		\end{tabular}
	\end{center}
	\choiceTF 
		{Cỡ mẫu  của mẫu số liệu là $ 12 $}
		{\True Nhóm chứa mốt của mẫu số liệu này là  $[40;60)$}
		{Nhóm chứa trung vị  là $[60;80)$}
		{Tứ phân vị  thứ nhất là $ 32 $}
	\loigiai{ 
		\begin{itemchoice}
			\itemch \textbf{Sai. }\\ 
			Cỡ  mẫu $n=5+9+12+10+6=42$.
			\itemch \textbf{Đúng.  }\\
			Tần số lớn nhất là $12$ nên  nhóm chứa mốt là $[40;60)$.
			 \itemch \textbf{Sai. }\\
			Gọi $x_1,\ldots, x_{42}$ là thời  gian tập thể dục trong ngày của $ 42 $ học sinh và giả sử dãy này được sắp xếp theo thứ tự tăng dần. Khi đó, trung vị là $\dfrac{x_{21}+x_{22}}{2}$. Do $ 2 $ giá trị $x_{21}, x_{22}$ thuộc nhóm $[40;60)$ nên nhóm $[40;60)$ chứa trung vị.
			 \itemch \textbf{Sai. }\\
			Tứ phân vị thứ nhất  $Q_1$ là $x_{11}$ nên nhóm chứa $Q_1$ là $[20;40)$.\\
			Do  đó $$Q_1=20+\dfrac{\dfrac{42}{4}-5}{9}\cdot 20=\dfrac{290}{9}.$$
		\end{itemchoice}
	} 
\end{ex}

% nguồn: Ok_dungsai_lop11_Lan2
\begin{ex}%[1D5V1-1]%[ma:2b4f1280]
	Số câu trả lời đúng một bài thi trắc nghiệm môn Sinh học gồm $ 50 $ câu của lớp $11\; A$ ở một trường THPT như sau: 
	\begin{center}
		\begin{tabular}{|c|c|c|c|c|c|}
			\hline 
			Số câu đúng & $[14;21)$ & $[21;28)$ & $[28;35)$  & $[35;42)$ & $[42;49)$ \\
			 \hline
			Số học sinh &  $4$ & $8$ & $25$ & $6$ & $7$ \\
			 \hline
		\end{tabular}
	\end{center}
	 \choiceTF
		{\True  Giá trị đại diện của nhóm $[14;21)$ là $ 17{,}5 $}
		{\True Giá trị đại diện  của nhóm $[21;28)$ là $ 24{,}5 $}
		{\True Giá trị đại diện  của nhóm $[42;49)$ là $ 45{,}5 $}
		{Số câu đúng trung bình là  $32{,}26$}
	\loigiai{ 
		\begin{center}
			\begin{tabular}{|c|c|c|c|c|c|}
				\hline 
				Số câu đúng & {$[14;21)$} & $[21;28)$ & $[28;35)$ & {$[35;42)$} &  $[42;49)$ \\
				\hline 
				Giá trị đại diện & $ {17{,}5} $ & $ 24{,}5 $ & $ 31{,}5 $ & $ 38{,}5 $ & $  45{,}5 $ \\
				 \hline
				 Số học sinh & {$ 4 $} & $ 8 $ &$ 25 $ & $ 6 $ & $ 7 $\\
				\hline 
			\end{tabular}
		\end{center}
		\begin{itemchoice}
			\itemch \textbf{Đúng.} Giá trị đại diện của nhóm $[14;21)$ là  $ 17{,}5 $.
			\itemch \textbf{Đúng.} Giá trị đại diện của nhóm $[21;28)$ là  $ 24{,}5 $.
			\itemch \textbf{Đúng.} Giá trị đại diện của  nhóm $[42;49)$ là $ 45{,}5 $.
			\itemch \textbf{Sai.}  \\
			 Số câu đúng trung bình
			\[\overline{x}=\dfrac{4\cdot17{,}5+8\cdot24{,}5+25\cdot31{,}5+6\cdot38{,}5+7\cdot45{,}5}{50}=32{,}06. 
			 \]
		\end{itemchoice}
	} 
\end{ex}

% nguồn: Ok_dungsai_lop11_Lan2
\begin{ex}%[1D5V1-1]%[ma:c473dfc9]
	Số lượng người đi xem một bộ phim mới theo độ tuổi trong một rạp chiếu phim (sau $1\; $giờ đầu công chiếu) được ghi lại theo bảng phân  phối ghép nhóm sau
	\begin{center}
		\begin{tabular}{|c|c|c|c|c|c|}
			 \hline
			Độ  tuổi & $[10;20)$ & $[20;30)$ & $[30;40)$ & $[40;50)$ & $[50;60)$ \\
			 \hline
			 Số người & $ 6 $ & $ 12 $ & $ 16 $ & $ 7 $ & $ 2 $ \\
			 \hline
		\end{tabular}
	\end{center}
	\choiceTF 
		{\True  Giá trị đại diện nhóm $[50;60)$ là $55$}
		{\True Độ tuổi được dự báo là ít xem phim  đó nhất là thuộc nhóm $[50;60)$}
		{\True Nhóm chứa mốt là nửa  khoảng $[30;40)$}
		 {Độ tuổi được dự báo là thích xem phim đó nhiều nhất là 31 tuổi}
	\loigiai{ 
		\begin{itemchoice}
			\itemch \textbf{Đúng.}\\  Giá trị đại diện nhóm $[50;60)$ là $55$.
			\itemch \textbf{Đúng.}\\ Độ tuổi được dự báo là ít xem phim đó nhất là $ 2 $ thuộc  nhóm $[50;60)$.
			\itemch \textbf{Đúng.}\\  Nhóm chứa mốt là nửa khoảng $[30;40)$.
			\itemch \textbf{Sai.}\\ Ta có\\ $u_m=30 $ , $n_m=16$ , $n_{m-1}=12$ , $n_{m+1}=7$  , $u_{m+1}-u_m=40-30=10$.\\
			Ta có mốt là $$M_0=30+\dfrac{16-12}{(16-12)+(16-7)} \cdot 10=\dfrac{430}{13}  \approx 33{,}08.$$
			Vậy độ tuổi được dự báo là thích xem phim đó nhiều nhất là $ 33  $ tuổi.
		\end{itemchoice}
	} 
\end{ex}


%%%%%%%%%%%%%%%% TỰ LUẬN %%%%%%%%%%%%%%%%

% nguồn: Ok_TuanLuan_Toan11_Lan1
\begin{bt}%[1D5N1-1]%[ma:20ce4f72]
	Cho  mẫu số liệu ghép nhóm về thống kê nhiệt độ tại một địa điểm trong $40$ ngày, ta có bảng số liệu sau:
	\begin{center}
		\begin{tabular}{|c|c|c|c|c|}
			\hline 
			Nhiệt độ $(^\circ C)$ & $[19{;}22$ & $[22{;}25)$ & $[25{;}28)$ & $[28{;}31)$  \\
			 \hline
			Số ngày & $7$  & $15$ & $12$ & $6$ \\
			\hline 
		\end{tabular}
	\end{center}
	Có bao nhiêu ngày có nhiệt độ từ $28^\circ  C$ đến dưới $31^\circ C$?
	 \loigiai{
		Có  6 ngày có nhiệt độ từ $28^\circ C$ đến dưới $31^\circ C$.
	 }
\end{bt}

% nguồn: Ok_TuanLuan_Toan11_Lan2
\begin{ex}%[1D5N1-1]%[ma:088ed9a2]
	Cho mẫu số liệu ghép nhóm  về thống kê nhiệt độ tại một địa điểm trong $40$ ngày, ta có bảng số liệu sau:
	\begin{center}
		\begin{tabular}{|c|c|c|c|c|}
			\hline 
			Nhiệt  độ $(^\circ C)$ & $[19{;}22$ & $[22{;}25)$ & $[25{;}28)$ & $[28{;}31)$ \\
			 \hline
			Số ngày & $7$  & $15$ & $12$ & $6$ \\
			 \hline
		\end{tabular}
	\end{center}
	Có bao nhiêu ngày  có nhiệt độ từ $28^\circ C$ đến dưới $31^\circ C$?
	\loigiai{ 
		Có 6 ngày có nhiệt độ  từ $28^\circ C$ đến dưới $31^\circ C$.
	} 
\end{ex}

% nguồn: Ok_TuanLuan_Toan11_Lan1
\begin{bt}%[1D5H1-1]%[ma:e028e5bd]
	Tuổi thọ của $30$ bóng đèn được thắp thử (đơn vị: giờ) được cho bởi bảng số liệu  thống kê dưới đây
	\begin{center}
		\begin{tabular}{|c|c|c|c|c|c|c|c|c|c|}
			 \hline
			$1180$ & $1150$ & $1190$ & $1170$ & $1180$ & $1170$ & $1160$ &  $1170$ & $1160$ & $1150$ \\
			 \hline
			$1190$ & $1180$ & $1170$ & $1170$ & $1170$ &  $1190$ & $1170$ & $1170$ & $1170$ & $1180$ \\
			\hline 
			$1170$ & $1160$ & $1160$  & $1160$ & $1170$ & $1160$ & $1180$ & $1180$ & $1150$ & $1170$ \\
			\hline 
		\end{tabular}
	\end{center}
	 \loigiai{
		Từ bảng số liệu trên ta suy ra bảng phân bố tần số  các giá trị tuổi thọ của 30 bóng đèn như sau
		\begin{center}
			\begin{tabular}{|c|c|c|c|c|c|c|}
				 \hline
				Tuổi thọ (giờ) & $1150$ &  $1160$ & $1170$ & $1180$ & $1190$ & Tổng \\
				 \hline
				 Tần số & $3$ & $6$ & $12$ & $6$ & $3$ & $30$ \\
				 \hline
			\end{tabular}
		\end{center}
		Ta thấy giá trị $1170$ có tần số bằng $12$ là lớn  nhất. \\
		Do đó  mốt của bảng số liệu là: $M_\circ=1170$.
	} 
\end{bt}

% nguồn: Ok_TuanLuan_Toan11_Lan1
\begin{bt}%[1D5H1-1]%[ma:dedea078]
	Thời gian (phút) để học sinh hoàn thành một câu  hỏi thi được cho như sau:
	\begin{center}
		\begin{tabular}{|c|c|c|c|c|c|}
			\hline 
			Thời  gian (phút) & $[0{,}5{;}10{,}5)$ & $[10{,}5{;}20{,}5)$ & $[20{,}5{;}30{,}5)$ & $[30{,}5{;}40{,}5)$ & $[40{,}5{;}50{,}5)$ \\
			 \hline
			 Số học sinh & $2$ & $10$ & $6$ & $4$ & $3$ \\
			 \hline
		\end{tabular}
	\end{center}
	Tính mốt của mẫu số liệu ghép nhóm này. 
	 \loigiai{
		Tần  số lớn nhất là $10$ nến nhóm chứa mốt là nhóm $[10{,}5{;}10{,}5)$. \\
		Ta có $j=2$;  $a_2=10{,}5$; $m_2=10$; $m_1=2$; $m_3=6$; $h=10$. \\
		Do đó  $M_\circ = 10{,}5+\dfrac{10-2}{(10-2)+(10-6)}\cdot10\approx17{,}17$.
	 }
\end{bt}

% nguồn: Ok_TuanLuan_Toan11_Lan1
\begin{bt}%[1D5H1-1]%[ma:7f2151b4]
	Một cửa hàng đã thống kê số ba  lô bán được mỗi ngày trong tháng $9$ với kết quả cho như sau:
	\begin{center}
		\begin{tabular}{|c|c|c|c|c|c|c|c|c|c|c|c|c|c|c|}
			 \hline
			$12$ & $29$ & $12$ & $19$ & $15$ & $21$ & $19$ & $29$ & $28$  & $12$ & $15$ & $25$ & $16$ & $20$ & $29$ \\
			\hline 
			$21$ & $12$ & $24$ & $14$ & $10$ & $12$ & $10$ & $23$ & $27$ & $28$ & $18$ &  $16$ & $10$ & $20$ & $21$ \\
			\hline 
		\end{tabular}
	\end{center}
	Hãy chia mẫu số liệu trên thành $5$ nhóm, lập bảng tần số ghép nhóm, hiệu chỉnh bảng tần số ghép nhóm  và xác định giá trị đại diện cho mỗi nhóm.
	\loigiai{ 
		Khoảng biến thiên của  mẫu số liệu trên là $R=29-10=19$. \\
		Độ dài mỗi nhóm $L>\dfrac{R}{k}=\dfrac{19}{5}=3{,}8$. \\ 
		Ta chon $L=4$ và  chia dữ liệu thành các nhóm: $[10{;}14)$, $[14{;}18)$, $[18{;}22)$, $[22{;}26)$, $[26{;}30)$. \\
		Khi  đó ta có bảng tần số ghép nhóm sau:
		\begin{center}
			\begin{tabular}{|c|c|c|c|c|c|}
				 \hline
				Số ba lô & $[10{;}14)$ & $[14{;}18)$  & $[18{;}22)$ & $[22{;}26)$ & $[26{;}30)$ \\
				\hline 
				Giá trị đại  diện & $12$ & $16$ & $20$ & $24$ & $28$ \\
				 \hline
				Số ngày & $8$ & $5$ & $8$ & $3$  & $6$ \\
				\hline 
			\end{tabular}
		\end{center}
	 }
\end{bt}

% nguồn: Ok_TuanLuan_Toan11_Lan1
\begin{bt}%[1D5H1-1]%[ma:e24a5194]
	Kết quả kiểm tra  chất lượng đầu năm (thang điểm $30$) của $41$ học sinh của một lớp được cho bởi bảng số liệu thống kê dưới đây
	\begin{center}
		\begin{tabular}{|c|c|c|c|c|c|c|c|c|c|c|c|c|}
			\hline 
			Điểm & $9$ & $11$ & $14$ & $16$ & $17$ & $18$ & $20$  & $21$ & $23$ & $25$ & Tổng \\
			\hline 
			Tần số & $3$ &  $7$ & $4$ & $4$ & $6$ & $7$ & $3$ & $3$ & $2$ & $2$ & $41$ \\
			 \hline
		\end{tabular}
	\end{center}
	Hãy tính mốt của bảng số  liệu thống kê trên.
	\loigiai{ 
		Ta thấy điểm $11$ và $18$ có tần số bằng  $7$ là lớn nhất. \\
		Do đó  bảng số liệu có hai mốt là: $M_{\circ}^{(1)}=11$ và $M_{\circ}^{(2)}=18$.
	} 
\end{bt}

% nguồn: Ok_TuanLuan_Toan11_Lan1
\begin{bt}%[1D5H1-1]%[ma:1587203e]
	Một bác sĩ mắt ghi lại tuổi của 30 bệnh nhân mắc bệnh đau mắt hột. Kết quả thu được mẫu số liệu  như sau:
	\begin{center}
		\begin{tabular}{|c|c|c|c|c|c|c|c|c|c|}
			\hline 
			$21$ & $17$ & $22$ & $18$ & $20$ & $17$ & $15$ & $13$ & $15$ &  $20$ \\
			 \hline
			$15$ & $12$ & $18$ & $17$ & $25$ &  $17$ & $21$ & $15$ & $12$ & $18$ \\
			\hline 
			$16$ & $23$ & $14$ &  $18$ & $19$ & $13$ & $16$ & $19$ & $18$ & $17$ \\
			 \hline
		\end{tabular}
	\end{center}
	Tính mốt $M_\circ$ của  bảng số liệu đã cho.
	 \loigiai{
		Từ bảng số liệu trên ta suy ra bảng phân bố tần số tuổi của 30 bệnh nhân đau  mắt hột như sau
		\begin{center}
			\begin{tabular}{|c|c|c|c|c|c|c|c|c|c|c|c|c|c|c|}
				 \hline
				Tuổi & $12$  & $13$ & $14$ & $15$ & $16$ & $17$ & $18$ & $19$ & $20$ & $21$ & $22$ & $23$ & $24$ & $25$ \\
				 \hline
				Tần số & $2$ & $2$ & $1$ & $4$ & $2$ & $5$ & $5$  & $2$ & $2$ & $2$ & $1$ & $1$ & $1$ & $30$ \\
				\hline 
			\end{tabular}
		\end{center}
		Ta thấy tuổi $17$ và $18$ có tần  số bằng $5$ lớn nhất. Do đó bảng số liệu có hai mốt là $M_{\circ}^{(1)}=17$ và $M_{\circ}^{(2)}=18$.
	 }
\end{bt}

% nguồn: Ok_TuanLuan_Toan11_Lan2
\begin{ex}%[1D5H1-1]%[ma:2456c824]
	Tuổi thọ của $30$ bóng đèn được thắp thử (đơn vị: giờ) được cho bởi bảng  số liệu thống kê dưới đây
	\begin{center}
		\begin{tabular}{|c|c|c|c|c|c|c|c|c|c|}
			\hline 
			$1180$ & $1150$  & $1190$ & $1170$ & $1180$ & $1170$ & $1160$ & $1170$ & $1160$ & $1150$ \\
			 \hline
			$1190$ & $1180$ & $1170$ & $1170$ &  $1170$ & $1190$ & $1170$ & $1170$ & $1170$ & $1180$ \\
			\hline 
			$1170$ & $1160$ & $1160$ & $1160$ & $1170$ & $1160$ &  $1180$ & $1180$ & $1150$ & $1170$ \\
			 \hline
		\end{tabular}
	\end{center}
	 \loigiai{
		Từ bảng số  liệu trên ta suy ra bảng phân bố tần số các giá trị tuổi thọ của 30 bóng đèn như sau
		\begin{center}
			\begin{tabular}{|c|c|c|c|c|c|c|}
				 \hline
				Tuổi thọ (giờ) & $1150$ & $1160$ & $1170$ & $1180$ & $1190$ &  Tổng \\
				 \hline
				Tần số & $3$ & $6$ & $12$ & $6$ &  $3$ & $30$ \\
				 \hline
			\end{tabular}
		\end{center}
		Ta thấy giá trị $1170$ có tần  số bằng $12$ là lớn nhất. \\
		Do  đó mốt của bảng số liệu là: $M_\circ=1170$.
	} 
\end{ex}

% nguồn: Ok_TuanLuan_Toan11_Lan2
\begin{ex}%[1D5H1-1]%[ma:2a0ec460]
	Thời gian (phút)  để học sinh hoàn thành một câu hỏi thi được cho như sau:
	\begin{center}
		\begin{tabular}{|c|c|c|c|c|c|}
			\hline 
			Thời gian (phút)  & $[0{,}5{;}10{,}5)$ & $[10{,}5{;}20{,}5)$ & $[20{,}5{;}30{,}5)$ & $[30{,}5{;}40{,}5)$ & $[40{,}5{;}50{,}5)$ \\
			 \hline
			Số học sinh & $2$ & $10$ & $6$ & $4$ &  $3$ \\
			\hline 
		\end{tabular}
	\end{center}
	Tính mốt của mẫu số  liệu ghép nhóm này.
	 \loigiai{
		Tần số lớn nhất là $10$ nến nhóm chứa mốt là nhóm $[10{,}5{;}10{,}5)$.  \\
		Ta có  $j=2$; $a_2=10{,}5$; $m_2=10$; $m_1=2$; $m_3=6$; $h=10$. \\
		 Do đó $M_\circ = 10{,}5+\dfrac{10-2}{(10-2)+(10-6)}\cdot10\approx17{,}17$.
	 }
\end{ex}

% nguồn: Ok_TuanLuan_Toan11_Lan2
\begin{ex}%[1D5H1-1]%[ma:c0d62984]
	 Một cửa hàng đã thống kê số ba lô bán được mỗi ngày trong tháng $9$ với kết quả cho như sau:
	\begin{center}
		\begin{tabular}{|c|c|c|c|c|c|c|c|c|c|c|c|c|c|c|}
			 \hline
			$12$ & $29$  & $12$ & $19$ & $15$ & $21$ & $19$ & $29$ & $28$ & $12$ & $15$ & $25$ & $16$ & $20$ & $29$ \\
			 \hline
			$21$ &  $12$ & $24$ & $14$ & $10$ & $12$ & $10$ & $23$ & $27$ & $28$ & $18$ & $16$ & $10$ & $20$ & $21$ \\
			\hline 
		\end{tabular}
	\end{center}
	Hãy chia mẫu  số liệu trên thành $5$ nhóm, lập bảng tần số ghép nhóm, hiệu chỉnh bảng tần số ghép nhóm và xác định giá trị đại diện cho mỗi nhóm.
	 \loigiai{
		Khoảng biến thiên  của mẫu số liệu trên là $R=29-10=19$. \\
		 Độ dài mỗi nhóm $L>\dfrac{R}{k}=\dfrac{19}{5}=3{,}8$. \\
		Ta chon $L=4$ và chia dữ liệu thành các nhóm: $[10{;}14)$, $[14{;}18)$,  $[18{;}22)$, $[22{;}26)$, $[26{;}30)$. \\
		Khi đó  ta có bảng tần số ghép nhóm sau:
		\begin{center}
			\begin{tabular}{|c|c|c|c|c|c|}
				 \hline
				Số ba lô & $[10{;}14)$ & $[14{;}18)$  & $[18{;}22)$ & $[22{;}26)$ & $[26{;}30)$ \\
				 \hline
				Giá trị đại diện & $12$ &  $16$ & $20$ & $24$ & $28$ \\
				 \hline
				Số ngày & $8$ & $5$ & $8$ &  $3$ & $6$ \\
				 \hline
			\end{tabular}
		\end{center}
	} 
\end{ex}

% nguồn: Ok_TuanLuan_Toan11_Lan2
\begin{ex}%[1D5H1-1]%[ma:ebc15312]
	Kết quả kiểm tra chất lượng đầu năm (thang điểm $30$) của $41$ học sinh của một lớp  được cho bởi bảng số liệu thống kê dưới đây
	\begin{center}
		\begin{tabular}{|c|c|c|c|c|c|c|c|c|c|c|c|c|}
			\hline 
			Điểm & $9$ & $11$ & $14$ & $16$ & $17$ & $18$ & $20$ & $21$ & $23$ & $25$ &  Tổng \\
			 \hline
			Tần số & $3$ &  $7$ & $4$ & $4$ & $6$ & $7$ & $3$ & $3$ & $2$ & $2$ & $41$ \\
			\hline 
		\end{tabular}
	\end{center}
	Hãy tính mốt của bảng số liệu  thống kê trên.
	 \loigiai{
		Ta thấy điểm $11$ và $18$ có tần số bằng $7$ là lớn nhất. \\ 
		Do đó bảng số liệu có hai mốt là: $M_{\circ}^{(1)}=11$  và $M_{\circ}^{(2)}=18$.
	} 
\end{ex}

% nguồn: Ok_TuanLuan_Toan11_Lan2
\begin{ex}%[1D5H1-1]%[ma:46478bb8]
	Một bác sĩ mắt ghi lại tuổi của 30 bệnh  nhân mắc bệnh đau mắt hột. Kết quả thu được mẫu số liệu như sau:
	\begin{center}
		\begin{tabular}{|c|c|c|c|c|c|c|c|c|c|}
			 \hline
			$21$ & $17$ & $22$ & $18$ & $20$ & $17$  & $15$ & $13$ & $15$ & $20$ \\
			\hline 
			$15$ & $12$ & $18$ & $17$ & $25$ & $17$ & $21$ & $15$ & $12$ &  $18$ \\
			 \hline
			$16$ & $23$ & $14$ & $18$ & $19$ & $13$ & $16$ & $19$  & $18$ & $17$ \\
			\hline 
		\end{tabular}
	\end{center}
	Tính mốt $M_\circ$ của bảng số liệu đã  cho.
	\loigiai{ 
		Từ bảng số liệu trên ta suy ra bảng phân bố tần số tuổi của 30  bệnh nhân đau mắt hột như sau
		\begin{center}
			\begin{tabular}{|c|c|c|c|c|c|c|c|c|c|c|c|c|c|c|}
				 \hline
				Tuổi & $12$ & $13$ & $14$ &  $15$ & $16$ & $17$ & $18$ & $19$ & $20$ & $21$ & $22$ & $23$ & $24$ & $25$ \\
				\hline 
				Tần số & $2$ & $2$ & $1$ & $4$ & $2$ & $5$ & $5$ & $2$ & $2$ & $2$ & $1$ & $1$ & $1$ & $30$ \\ 
				 \hline
			\end{tabular}
		\end{center}
		Ta thấy tuổi $17$ và $18$ có  tần số bằng $5$ lớn nhất. Do đó bảng số liệu có hai mốt là $M_{\circ}^{(1)}=17$ và $M_{\circ}^{(2)}=18$.
	} 
\end{ex}

% nguồn: Ok_TuanLuan_Toan11_Lan1
\begin{bt}%[1D5B1-1]%[ma:8c565526]
	Tính đạo hàm bằng định nghĩa của hàm số $y=f(x)=x^3-3x^2$ tại  điểm $x_0=1$.
	\loigiai{Ta có $\displaystyle \lim \limits_{x \to 1}\dfrac{x^3-3x^2-(-2)}{x-1}=\displaystyle \lim \limits_{x \to 1}\dfrac{(x^2-2x-2)(x-1)}{x-1}=\displaystyle \lim \limits_{x \to 1}(x^2-2x-2)=-3$. 
	} 
\end{bt}

% nguồn: Ok_TuanLuan_Toan11_Lan1
\begin{bt}%[1D5B1-1]%[ma:a8473e5f]
	Dùng định nghĩa, tính đạo hàm  các hàm số sau đây:
	\begin{listEX}
		\item $y=x^2+3x-2$ tại $x_0=1$. 
		\item $y=-\dfrac{3}{x}$  tại $x_0=2$.
		\item $y=\dfrac{x+1}{x-1}$ tại  điểm $x_0=0$.
	\end{listEX}
	\loigiai{ 
		\begin{listEX}
			 \item
			Ta  có $\displaystyle \lim \limits_{x \to 1}\dfrac{x^2+3x-2-2}{x-1}=\displaystyle \lim \limits_{x \to 1}\dfrac{(x+4)(x-1)}{x-1}=\displaystyle \lim \limits_{x \to 1}(x+4)=5$.
			\item 
			Ta có $\displaystyle \lim \limits_{x \to 2}\dfrac{\frac{-3}{x}+\frac{3}{2}}{x-2}=\displaystyle \lim \limits_{x \to  2}\dfrac{3x-6}{2x(x-2)}=\displaystyle \lim \limits_{x \to 2}\dfrac{3}{2x}=\dfrac{3}{4}$.
			\item Ta có  $\displaystyle \lim \limits_{x \to 0}\dfrac{\frac{x+1}{x-1}+1}{x}=\displaystyle \lim \limits_{x \to 0}\dfrac{2x}{x(x-1)}=-2$.
		\end{listEX}
	} 
\end{bt}

% nguồn: Ok_TuanLuan_Toan11_Lan2
\begin{ex}%[1D5B1-1]%[ma:38b6fe79]
	Tính đạo hàm bằng định  nghĩa của hàm số $y=f(x)=x^3-3x^2$ tại điểm $x_0=1$.
	\loigiai{Ta có $\displaystyle \lim \limits_{x \to 1}\dfrac{x^3-3x^2-(-2)}{x-1}=\displaystyle \lim \limits_{x \to 1}\dfrac{(x^2-2x-2)(x-1)}{x-1}=\displaystyle \lim \limits_{x \to 1}(x^2-2x-2)=-3$. 
	} 
\end{ex}

% nguồn: Ok_TuanLuan_Toan11_Lan2
\begin{ex}%[1D5B1-1]%[ma:45b3901d]
	Dùng định  nghĩa, tính đạo hàm các hàm số sau đây:
	\begin{listEX}
		\item $y=x^2+3x-2$  tại $x_0=1$.
		\item $y=-\dfrac{3}{x}$  tại $x_0=2$.
		\item $y=\dfrac{x+1}{x-1}$ tại điểm  $x_0=0$.
	\end{listEX}
	 \loigiai{
		\begin{listEX}
			 \item
			Ta có $\displaystyle \lim \limits_{x \to 1}\dfrac{x^2+3x-2-2}{x-1}=\displaystyle \lim \limits_{x \to 1}\dfrac{(x+4)(x-1)}{x-1}=\displaystyle \lim \limits_{x \to  1}(x+4)=5$.
			 \item
			Ta có $\displaystyle \lim \limits_{x \to 2}\dfrac{\frac{-3}{x}+\frac{3}{2}}{x-2}=\displaystyle \lim \limits_{x \to 2}\dfrac{3x-6}{2x(x-2)}=\displaystyle \lim \limits_{x \to 2}\dfrac{3}{2x}=\dfrac{3}{4}$. 
			\item Ta có $\displaystyle \lim \limits_{x \to 0}\dfrac{\frac{x+1}{x-1}+1}{x}=\displaystyle \lim \limits_{x \to 0}\dfrac{2x}{x(x-1)}=-2$. 
		\end{listEX}
	} 
\end{ex}
